%% =========================================================
%% Chapter: The Multi-Block Asymmetric Fundamental Theorem
%% Covers Notes/OpenProblemsTN/problems/p5_asymmetric_fundamental_theorem.tex, \S7,
%% and the single-block corollary of
%% Notes/OpenProblemsTN/strategies/p5_asymmetric_compression_theorem.tex.
%% =========================================================

\chapter{The Multi-Block Asymmetric Fundamental Theorem}%
\label{ch:asymmetric_ft}

The asymmetric Fundamental Theorem of matrix product states and operators
compares an arbitrary tensor with a weighted direct sum of normal tensors.
Equality of all positive-length periodic vectors determines every nonzero
simple composition factor of the arbitrary tensor, including its multiplicity;
the only undetectable factor is the one-dimensional tensor on which every
physical matrix vanishes. The remaining freedom is an upper-triangular,
globally nilpotent extension between the factors, explaining the failure of
sitewise intertwiners under periodic equality alone
(Remark~\ref{rem:asym_zipper}).

\begin{center}
    % Formula: the periodic trace tr(B^w), generating the periodic
    % vector V^{(N)}(B) compared against the target family throughout
    % this chapter.
    % Source: chapter introduction; the same picture as
    % Definition~\ref{def:periodic_mps_vector} (Chapter~\ref{ch:mps}),
    % drawn here for the arbitrary tensor B of the compression theorem.
    % Ink-to-index map: each node is the same local tensor B; the upper
    % legs carry the letters i_1,...,i_N of the word w, and the
    % ellipsis and brace make the N-site ring schematic.
    % Contraction set: the horizontal virtual legs are contracted around
    % the ring, closing the trace; no leg is left open.
    % Boundary signature: (virtual-west=0 | virtual-east=0 |
    % physical-up=3 | physical-down=0) for the finite schematic shown.
    \tenkzkernel
    \begin{tenkz}[cols=4, boundary=periodic, physical=up]
        \tn[label pos=135, ports={90:physical:$i_1$}]{B} &
        \tn[label pos=135]{B} & \tn[skin=dots]{} &
        \tn[label pos=135, ports={90:physical:$i_N$}]{B}
        \tnmark[form=bracket]{(1,1) .. (1,4)}{$N\text{ sites}$}
    \end{tenkz}
\end{center}

\section{Triangular word products}%
\label{sec:asym_triangular}

The distinction between a word-level compression and a sitewise
intertwiner is already visible in elementary block matrix multiplication.

\begin{lemma}[Diagonal blocks of block-triangular products]
    \label{lem:asym_triangular}
    \lean{Matrix.blockDiag'_mul_of_blockTriangular, Matrix.blockTriangular_evalWord,
        Matrix.blockDiag'_evalWord, Matrix.blockProj_mul_blockEmbed_self,
        Matrix.blockProj_mul_blockEmbed_of_ne, Matrix.blockProj_mul_mul_blockEmbed}
    \leanok
    \uses{def:eval_word}
    Let $o$ be a finite index set with block sizes $n_k$, $k\in o$, and let
    $T^1,\ldots,T^d$ be block upper triangular for the resulting
    decomposition of $\MN{\sum_k n_k}$. Write $T^i_{kk}$ for the $k$-th
    diagonal block of $T^i$. Then $T^w$ is block upper triangular for every
    word $w$, and its $k$-th diagonal block is the word evaluation of the
    diagonal blocks,
    \begin{align}
        (T^w)_{kk} & =T^{i_1}_{kk}T^{i_2}_{kk}\cdots T^{i_N}_{kk}
        \label{eq:asym_triangular_word}
    \end{align}
    for $w=(i_1,\ldots,i_N)$. If $P_k$ is the coordinate embedding of the
    $k$-th block and $Q_k$ the coordinate projection onto it, then for every
    matrix $M$ on the same space, $Q_kMP_k=M_{kk}$, and
    \begin{align}
        Q_kP_l & =\begin{cases}\id, & k=l,\\0, & k\ne l.\end{cases}
        \label{eq:asym_block_biorthogonality}
    \end{align}
\end{lemma}

\begin{proof}
    \leanok
    For two block upper-triangular matrices $X,Y$,
    $(XY)_{kk}=\sum_l X_{kl}Y_{lk}$; the terms with $l<k$ vanish because
    $X_{kl}=0$, and those with $l>k$ vanish because $Y_{lk}=0$, leaving
    $(XY)_{kk}=X_{kk}Y_{kk}$. Induction on the word length gives
    (\ref{eq:asym_triangular_word}) and the persistence of block
    triangularity. The identity $Q_kMP_k=M_{kk}$ and
    (\ref{eq:asym_block_biorthogonality}) are direct computations from the
    coordinate description of $P_k$ and $Q_k=P_k^{\mathsf T}$.
\end{proof}

\begin{remark}[Periodic equality does not imply a sitewise intertwiner]
    \label{rem:asym_zipper}
    Take the one-dimensional normal tensor $A^1=(1)$, $A^2=(2)$, and the
    bond-dimension-three tensor
    \begin{align}
        B^1
         & =\begin{pmatrix}0&1&0\\0&1&0\\0&0&0\end{pmatrix},
         &
        B^2
         & =\begin{pmatrix}0&0&0\\0&2&1\\0&0&0\end{pmatrix}.
        \notag
    \end{align}
    Both matrices are upper triangular with diagonal blocks $0$, $A^i$, $0$.
    With $P=e_2$ and $Q=e_2^{\mathsf T}$, Lemma~\ref{lem:asym_triangular}
    gives $QB^wP=A^w$ and $QP=1$ for every nonempty word $w$, so
    $\tr(B^w)=A^w=\tr(A^w)$: the periodic vectors of $A$ and $B$ agree at
    every positive length.

    There is nevertheless no nonzero right sitewise intertwiner: a vector
    $v=(x,y,z)^{\mathsf T}$ with $B^1v=v$ and $B^2v=2v$ would need $x=y$,
    $z=0$ from the first equation and $x=0$ from the first coordinate of the
    second, forcing $v=0$. Symmetrically, no nonzero left intertwiner
    $uB^1=u$, $uB^2=2u$ exists: the first equation gives $u_1=u_3=0$ while
    the third coordinate of the second gives $u_2=2u_3=0$. Periodic equality
    can therefore force a multiplicative subquotient while admitting neither
    one-sided sitewise intertwiner, so it cannot imply a biorthogonal zipper
    pair. This is why the theorem below is stated as a compression
    $W_sB^wV_s=C_s^w$ rather than as an eigenvector relation.
\end{remark}

\section{The word algebra and normal tensors as simple modules}%
\label{sec:asym_word_algebra}

\begin{definition}[Word algebra and positive ideal]
    \label{def:asym_word_algebra}
    \lean{WordAlgebra, WordAlgebra.aug}
    \leanok
    Fix $d\geq1$. The \emph{word algebra} is the free unital associative
    algebra $F=\C\langle x_1,\ldots,x_d\rangle$, realized as the algebra of
    finite $\C$-linear combinations of words in the $d$ letters under
    concatenation, so that the words form a basis. The \emph{augmentation}
    $\mathrm{aug}:F\to\C$ is the algebra map sending every generator to
    zero; on a basis word $w$, $\mathrm{aug}(w)=1$ if $w$ is empty and $0$
    otherwise. The \emph{positive ideal} is the two-sided ideal
    $F_+=\ker(\mathrm{aug})$ of elements with zero constant term.
\end{definition}

\begin{definition}[Word-algebra module of a tensor]
    \label{def:asym_word_module}
    \lean{MPSTensor.WordModule}
    \leanok
    \uses{def:asym_word_algebra, def:mps_tensor, def:eval_word}
    A tensor $C=(C^i)_{i=0}^{d-1}$ on $\C^D$ defines an $F$-module structure
    on $\C^D$ through the representation $\rho_C(x_i)=C^i$, extended
    multiplicatively so that a word $w$ acts by $\rho_C(w)=C^w$.
\end{definition}

\begin{definition}[Word traces]
    \label{def:asym_word_traces}
    \lean{WordAlgebra.traceChar, WordAlgebra.traceWord}
    \leanok
    \uses{def:asym_word_algebra}
    For a finite-dimensional $F$-module $M$, the \emph{trace character} is
    the linear functional $p\mapsto\tr(\rho_M(p))$ on $F$, and the
    \emph{word trace} of a word $w$ is $\tr(\rho_M(w))$.
\end{definition}

\begin{lemma}[Word traces of a tensor]
    \label{lem:asym_word_traces_tensor}
    \lean{MPSTensor.traceWord_wordModule}
    \leanok
    \uses{def:asym_word_traces, def:asym_word_module}
    On the word-algebra module of a tensor $C$, the word trace of $w$ is
    the periodic coefficient $\tr(C^w)$.
\end{lemma}

\begin{proof}
    \leanok
    \uses{def:asym_word_module}
    The action of $w$ is conjugate to the matrix $C^w$ under the
    identification of the module with $\C^D$, and the trace of a linear
    map is the trace of its matrix.
\end{proof}

\begin{lemma}[Normal tensors define simple modules]
    \label{lem:asym_normal_simple}
    \lean{MPSTensor.isSimpleModule_wordModule_of_isNormal}
    \leanok
    \uses{def:asym_word_module, def:normal}
    If $C$ is normal on $\C^D$ with $D>0$, its word-algebra module is
    simple.
\end{lemma}

\begin{proof}
    \leanok
    Let $K\subseteq\C^D$ be a nonzero invariant subspace and $u\in K$
    nonzero. Normality gives a length $L$ for which the span of the
    length-$L$ word evaluations is all of $\MN{D}$; since $K$ is invariant
    under every such word and hence under their span, $Xu\in K$ for every
    $X\in\MN{D}$. Every vector is $Xu$ for a suitable $X$ because $u\neq0$,
    so $K=\C^D$.
\end{proof}

Multiplying every matrix of a tensor by a nonzero scalar does not change its
invariant subspaces, so a weighted block $\mu_kA_{(k)}$ built from a normal
tensor $A_{(k)}$ and $\mu_k\in\C^\times$ remains a simple $F$-module.

\section{Separating simple modules}%
\label{sec:asym_separation}

\begin{lemma}[Positive-word character lemma]
    \label{lem:asym_character}
    \notready
    \uses{def:asym_word_module}
    Let $U,V$ be finite-dimensional $F$-modules with
    $\tr(\rho_U(p))=\tr(\rho_V(p))$ for every $p\in F_+$. Then the
    semisimplifications of $U$ and $V$ contain every simple $F$-module with
    the same multiplicity, except possibly the one-dimensional module
    $\mathbf0$ on which every generator acts as zero: there are integers
    $a,b\geq0$ and a semisimple module $M$ without a $\mathbf0$ summand such
    that $U^{\mathrm{ss}}\cong M\oplus\mathbf0^{\oplus a}$ and
    $V^{\mathrm{ss}}\cong M\oplus\mathbf0^{\oplus b}$.
    % The source proof decomposes the finite-dimensional image algebra
    % modulo its Jacobson radical by Artin--Wedderburn and compares
    % traces of central idempotents.
\end{lemma}

\begin{lemma}[Positive-ideal separation of linear functionals]
    \label{lem:asym_ker_aug_separation}
    \lean{WordAlgebra.linearMap_eq_on_ker_aug}
    \leanok
    \uses{def:asym_word_algebra}
    Two $\C$-linear functionals on $F$ agreeing on every nonempty word
    agree on the positive ideal $F_+$.
\end{lemma}

\begin{proof}
    \leanok
    Their difference agrees with $(\chi_1(1)-\chi_2(1))\cdot\mathrm{aug}$ on
    every basis word, by induction on the monoid-algebra presentation, hence
    on all of $F$ by linearity; on $F_+$ the augmentation vanishes, so the
    two functionals agree there.
\end{proof}

\begin{lemma}[The zero module is the only simple module with trivial
        generator action]
    \label{lem:asym_zero_module}
    \lean{WordAlgebra.finrank_eq_one_of_isSimpleModule_of_forall_smul_eq_zero}
    \leanok
    \uses{def:asym_word_algebra}
    A finite-dimensional simple $F$-module on which every generator acts as
    zero is one-dimensional. It is the module $\mathbf0$ of
    Lemma~\ref{lem:asym_character}.
\end{lemma}

\begin{proof}
    \leanok
    If every generator acts as zero, $F$ acts through the augmentation, so
    the $F$-submodules and the $\C$-subspaces coincide; simplicity as a
    $\C$-vector space forces dimension one.
\end{proof}

\begin{theorem}[Burnside's theorem at module level]
    \label{thm:asym_burnside}
    \lean{WordAlgebra.exists_smul_eq_of_isSimpleModule}
    \leanok
    \uses{def:asym_word_module}
    Every $\C$-linear endomorphism of a finite-dimensional simple
    $F$-module $Q$ is the action of some element of $F$.
\end{theorem}

\begin{proof}
    \leanok
    By Schur's lemma over the algebraically closed field $\C$, the
    commutant $\End_F(Q)$ consists of scalars, so $Q$ is finite-dimensional
    over its commutant. The Jacobson density theorem then makes the action
    map $F\to\End_{\C}(Q)$ surjective.
\end{proof}

\begin{theorem}[Separating a simple module from a finite non-isomorphic
        family]
    \label{thm:asym_separation}
    \lean{WordAlgebra.exists_aug_eq_zero_smul_eq_self,
        WordAlgebra.exists_aug_eq_zero_smul_eq_self_smul_eq_zero,
        MPSTensor.exists_aug_eq_zero_of_forall_isEmpty_linearEquiv}
    \leanok
    \uses{def:asym_word_module, def:asym_word_algebra}
    Let $Q$ be a finite-dimensional simple $F$-module on which some
    generator acts nonzero. There is $u\in F_+$ acting as the identity on
    $Q$. If $T$ is a finite-dimensional simple module not isomorphic to
    $Q$, there is $p\in F_+$ acting as the identity on $Q$ and as zero on
    $T$. More generally, if $C_s$, $s\in S$, is a finite family of
    word-algebra modules of tensors, each simple and none isomorphic to
    $Q$, there is $p\in F_+$ acting as the identity on $Q$ and as zero on
    every $C_s$.
\end{theorem}

\begin{proof}
    \leanok
    \uses{thm:asym_burnside}
    For the first claim, Burnside's theorem realizes, for a basis
    $\{b_j\}$ of $Q$ and a nonzero vector $x_i\cdot q_0$, the maps sending
    $x_i\cdot q_0\mapsto b_j$ and $q\mapsto(\text{coordinate }j\text{ of
        }q)\cdot q_0$; assembling them into
    $u=\sum_j\alpha_jx_i\beta_j$ gives an element of $F_+$ acting as the
    identity on $Q$. For the pairwise separation, Schur's lemma shows every
    $F$-linear endomorphism of $Q\times T$ preserves the two
    non-isomorphic simple factors, so $Q\times T$ is finite-dimensional
    over its (scalar) commutant; Jacobson density on $Q\times T$ produces
    $a\in F$ acting as the identity on the first factor and zero on the
    second, and $p=au$ acts as the identity on $Q$ and as zero on $T$. The
    finite-family statement follows by induction on $S$, multiplying the
    separators for the successive members.
\end{proof}

\begin{lemma}[Trace additivity over an invariant submodule]
    \label{lem:asym_trace_additivity}
    \lean{WordAlgebra.traceWord_eq_add_quotient}
    \leanok
    \uses{def:asym_word_traces}
    For a finite-dimensional $F$-module $M$ and a submodule $K\leq M$, the
    word trace of $M$ is the sum of the word traces of $K$ and of $M/K$.
\end{lemma}

\begin{proof}
    \leanok
    This is the additivity of the ordinary trace of an endomorphism over an
    invariant subspace, applied to the action of each word on $M$.
\end{proof}

\begin{theorem}[Nilpotency from vanishing trace powers]
    \label{thm:asym_nilpotent_trace_powers}
    \lean{LinearMap.isNilpotent_of_forall_trace_pow_eq_zero}
    \leanok
    An endomorphism of a finite-dimensional complex vector space, all of
    whose positive powers have vanishing trace, is nilpotent.
\end{theorem}

\begin{proof}
    \leanok
    The characteristic polynomial is determined by the power sums of its
    roots (Newton's identities), so vanishing trace powers force it to
    equal $X^{\dim}$; the Cayley--Hamilton theorem then gives nilpotency.
\end{proof}

\begin{theorem}[Identification of a simple quotient with a target block]
    \label{thm:asym_identification}
    \lean{MPSTensor.exists_linearEquiv_wordModule_of_isSimpleModule_quotient}
    \leanok
    \uses{def:asym_word_traces, def:asym_word_module}
    Let $M$ be a finite-dimensional $F$-module whose word traces are, on
    every nonempty word, the sum of the word traces of a finite family of
    simple tensor modules $C_s$, $s\in S$. Let $K\leq M$ have simple
    quotient $M/K$ on which some generator acts nontrivially. Then $M/K$ is
    isomorphic to $C_s$ for some $s\in S$.
\end{theorem}

\begin{proof}
    \leanok
    \uses{thm:asym_separation, lem:asym_ker_aug_separation,
        thm:asym_nilpotent_trace_powers}
    If $M/K$ were isomorphic to no $C_s$, Theorem~\ref{thm:asym_separation}
    supplies $p\in F_+$ acting as the identity on $M/K$ and as zero on
    every $C_s$. Since the word trace of $M$ agrees with the sum of the
    word traces of the $C_s$ on every nonempty word, and $p\in F_+$,
    Lemma~\ref{lem:asym_ker_aug_separation} shows that every positive power
    of the action of $p$ on $M$ has vanishing trace, so this action is
    nilpotent by Theorem~\ref{thm:asym_nilpotent_trace_powers}. A nilpotent
    action of $p$ on $M$ descends to a nilpotent action on the quotient
    $M/K$, contradicting that $p$ acts there as the identity on a nonzero
    module.
\end{proof}

\section{The composition-series flag}%
\label{sec:asym_flag}

\begin{lemma}[Composition-series triangular gauge]
    \label{lem:asym_composition_gauge}
    \lean{exists_linearEquiv_blockTriangular_of_flag}
    \leanok
    \uses{def:asym_word_module}
    Let $B=(B^i)_i$ act on $\C^{D_B}$, and choose a composition series
    $0=H_0\subset H_1\subset\cdots\subset H_r=\C^{D_B}$ of its
    word-algebra module. There is $G\in\GL_{D_B}(\C)$ such that every
    $T^i=GB^iG^{-1}$ is block upper triangular for the resulting
    decomposition into $r$ blocks, and the $k$-th diagonal block of $T^i$
    is the action induced by $B^i$ on the simple subquotient
    $H_k/H_{k-1}$.
\end{lemma}

\begin{proof}
    \leanok
    Choose a basis of $H_1$ and extend it successively to bases of
    $H_2,\ldots,H_r$. In the resulting basis, every $H_k$ is a coordinate
    initial segment; since every $H_k$ is invariant under every $B^i$, all
    matrices $B^i$ are block upper triangular in this basis, with $k$-th
    diagonal block the induced action on $H_k/H_{k-1}$.
\end{proof}

\begin{definition}[Labelled invariant flag]
    \label{def:asym_flag_data}
    \lean{MPSTensor.FlagData}
    \leanok
    \uses{def:asym_word_module}
    Let $M$ be a finite-dimensional $F$-module, $S$ a finite set of target
    slots, and $C_s$, $s\in S$, a family of tensors. A \emph{labelled
        invariant flag} of $M$ over $(S,C)$ is an invariant chain
    $0=H_0\leq H_1\leq\cdots\leq H_r=M$ together with a bijective labelling
    of the steps $\{1,\ldots,r\}$ by $S\sqcup\{1,\ldots,z\}$, for some
    $z\geq0$, such that a step labelled by $s\in S$ has subquotient
    $H_k/H_{k-1}$ isomorphic to the word-algebra module of $C_s$, and a
    step labelled by an index in $\{1,\ldots,z\}$ has one-dimensional
    subquotient on which every generator acts as zero.
\end{definition}

\begin{lemma}[A family of normal blocks with zero total word trace is
        empty]
    \label{lem:asym_flag_base_case}
    \lean{MPSTensor.finset_eq_empty_of_forall_sum_trace_evalWord_eq_zero}
    \leanok
    \uses{def:normal, def:asym_word_traces}
    If the word traces of a finite family of normal tensors of positive
    bond dimension sum to zero on every nonempty word, the family is
    empty.
\end{lemma}

\begin{proof}
    \leanok
    \uses{thm:asym_separation, thm:asym_nilpotent_trace_powers}
    Otherwise Theorem~\ref{thm:asym_separation} applied to one member $C_{s_0}$
    of the family supplies $u\in F_+$ acting as the identity on $C_{s_0}$;
    vanishing total word trace forces every positive power of the action of
    $u$ on the product of the modules in the family to have trace zero, so
    this action is nilpotent by Theorem~\ref{thm:asym_nilpotent_trace_powers}.
    A nilpotent action of $u$ contradicts its action as the identity on the
    nonzero summand $C_{s_0}$.
\end{proof}

\begin{theorem}[The flag theorem]
    \label{thm:asym_flag_theorem}
    \lean{MPSTensor.exists_flagData}
    \leanok
    \uses{def:asym_flag_data, def:normal, def:asym_word_traces}
    Let $C_s$, $s\in S$, be a finite family of normal tensors of positive
    bond dimension, and let $M$ be a finite-dimensional $F$-module whose
    word traces agree, on every nonempty word, with the sum of the word
    traces of the $C_s$. Then $M$ carries a labelled invariant flag over
    $(S,C)$.
\end{theorem}

\begin{proof}
    \leanok
    \uses{lem:asym_flag_base_case, lem:asym_zero_module,
        thm:asym_identification, lem:asym_trace_additivity}
    Induction on $\dim M$. If $M=0$, Lemma~\ref{lem:asym_flag_base_case}
    forces $S=\emptyset$, and the empty chain is the required flag. If
    $M\neq0$, choose a maximal proper invariant submodule $K\leq M$ (a
    coatom of the invariant-submodule lattice, which exists because $M$ is
    finite-dimensional), so $M/K$ is simple. If some generator acts
    nontrivially on $M/K$, Theorem~\ref{thm:asym_identification} identifies
    $M/K$ with some block $C_{s_0}$; Lemma~\ref{lem:asym_trace_additivity}
    shows the word trace of $K$ agrees with the sum over the remaining
    blocks $S\setminus\{s_0\}$, and the induction hypothesis (applied to
    $K$, of strictly smaller dimension) supplies a flag of $K$, extended by
    the matched final step $K\subset M$ labelled $s_0$. If instead every
    generator acts as zero on $M/K$, Lemma~\ref{lem:asym_zero_module} shows
    $M/K$ is one-dimensional; Lemma~\ref{lem:asym_trace_additivity} shows
    the word trace of $K$ already agrees with the sum over all of $S$, and
    the induction hypothesis supplies a flag of $K$, extended by the
    unmatched final step $K\subset M$.
\end{proof}

\section{The multi-block asymmetric compression theorem}%
\label{sec:asym_main}

\begin{theorem}[Multi-block asymmetric compression theorem]
    \label{thm:asym_multiblock}
    \lean{MPSTensor.MultiBlockCompression,
        MPSTensor.exists_multiBlockCompression_of_isNormal,
        MPSTensor.MultiBlockCompression.dim_eq,
        MPSTensor.MultiBlockCompression.isReduction,
        MPSTensor.MultiBlockCompression.left_mul_right_of_ne,
        MPSTensor.MultiBlockCompression.evalWord_remainder_eq_zero}
    \leanok
    \uses{def:same_mpv2_pos, def:normal}
    Let $A=\bigoplus_{s\in S}C_s$ be a finite direct sum of nonzero normal
    tensors $C_s=(C_s^i)_i$ of bond dimension $D_s$, with repetitions
    allowed, and let $B=(B^i)_i$ be an arbitrary tensor of bond dimension
    $D_B$. Assume that the periodic vectors of $A$ and $B$ agree at every
    positive length,
    $\ket{V^{(N)}(B)}=\ket{V^{(N)}(A)}$ for every $N>0$; equivalently,
    $\tr(B^w)=\tr(A^w)$ for every nonempty word $w$. Then there are an
    integer $z\geq0$, an integer $r=|S|+z$, an invertible matrix
    $G\in\GL_{D_B}(\C)$, and maps $V_s:\C^{D_s}\to\C^{D_B}$,
    $W_s:\C^{D_B}\to\C^{D_s}$, with the following properties.
    % Plain enumerate with explicit tags: the kernel key-value list syntax
    % [label=(\roman*)] requires LaTeX2e >= 2025, which CI does not have.
    \begin{enumerate}
        \item[(i)] Every $GB^iG^{-1}$ is block upper triangular for one common
              decomposition into $r$ blocks.
        \item[(ii)] Exactly $|S|$ of the diagonal slots are matched with the
              target slots: writing $j(s)$ for the position matched with $s$,
              there is $X_s\in\GL_{D_s}(\C)$ with
              $\bigl(GB^iG^{-1}\bigr)_{j(s),j(s)}=X_sC_s^iX_s^{-1}$.
        \item[(iii)] The remaining $z$ diagonal slots are one-dimensional zero
              tensors.
        \item[(iv)] The compression pairs are mutually biorthogonal: $W_sV_t=
                  \id_{D_s}$ if $s=t$ and $0$ otherwise.
        \item[(v)] For every nonempty word $w$, $W_sB^wV_s=C_s^w$.
        \item[(vi)] Writing $\widehat C^i=\sum_{s\in S}V_sC_s^iW_s$ and
              $R^i=B^i-\widehat C^i$, every product
              $R^{i_1}R^{i_2}\cdots R^{i_r}=0$: the non-unital algebra generated
              by the $R^i$ is nilpotent.
        \item[(vii)] $D_B=\sum_{s\in S}D_s+z$.
    \end{enumerate}
\end{theorem}

\begin{center}
    % Formula: clause (v), W_sB^wV_s=C_s^w for every nonempty word w.
    % Source: clauses (iv)-(v) of Theorem~\ref{thm:asym_multiblock}.
    % Ink-to-index map: the two end rings are the compression maps $W_s$
    % (left) and $V_s$ (right); the interior nodes are copies of $B$,
    % with upper legs carrying the letters of $w$.
    % Contraction set: every horizontal virtual bond between consecutive
    % copies of $B$ and between each $B$ and its adjacent ring is
    % contracted; the physical legs and the two outer virtual legs (the
    % $D_s$-dimensional indices of the matrix $C_s^w$) remain open.
    % Boundary signature: (virtual-west=1 | virtual-east=1 |
    % physical-up=2 | physical-down=0) for the finite schematic shown,
    % with the two explicit lettered nodes carrying the drawn physical
    % ports.
    \tenkzkernel
    \begin{tenkz}[rows={wire}, cols=5, west=open, east=open]
        \tn[skin=ring]{W_s} &
        \tn[label pos=s, ports={90:physical:$i_1$}]{B} &
        \tn[skin=dots]{} &
        \tn[label pos=s, ports={90:physical:$i_N$}]{B} &
        \tn[skin=ring]{V_s}
    \end{tenkz}
    \quad $=$ \quad
    \begin{tenkz}[rows={wire}, cols=3, west=open, east=open]
        \tn[label pos=s, ports={90:physical:$i_1$}]{C_s} &
        \tn[skin=dots]{} &
        \tn[label pos=s, ports={90:physical:$i_N$}]{C_s}
    \end{tenkz}
    \\[2ex]
    % Formula: clause (iv), W_sV_t=\id_{D_s} if s=t and 0 otherwise.
    % Source: clause (iv) of Theorem~\ref{thm:asym_multiblock}.
    % Ink-to-index map: the two rings are the compression maps $W_s$ and
    % $V_t$, joined with no intervening tensor; the two outer virtual
    % legs are the $D_s$- and $D_t$-dimensional indices of the matrix
    % $W_sV_t$.
    % Contraction set: the single virtual bond between the two rings.
    % Boundary signature: (virtual-west=1 | virtual-east=1 |
    % physical-up=0 | physical-down=0).
    \begin{tenkz}[rows={wire}, cols=2, west=open, east=open]
        \tn[skin=ring]{W_s} & \tn[skin=ring]{V_t}
    \end{tenkz}
    \quad $=\;\delta_{st}\,\id$
\end{center}

\begin{proof}
    \leanok
    \uses{thm:asym_flag_theorem, lem:asym_composition_gauge,
        lem:asym_triangular}
    By Theorem~\ref{thm:asym_flag_theorem}, applied with $M$ the
    word-algebra module of $B$, whose word traces agree with the sum of
    those of the $C_s$ by hypothesis, $\C^{D_B}$ carries a labelled
    invariant flag over $(S,C)$: an invariant chain $0=H_0\leq\cdots\leq
        H_r=\C^{D_B}$ whose subquotients are, in some order, the $C_s$ (each
    once) and $z=r-|S|$ copies of the zero module. This chain is in
    particular a composition series, so Lemma~\ref{lem:asym_composition_gauge}
    supplies $G\in\GL_{D_B}(\C)$ for which $T^i=GB^iG^{-1}$ is block upper
    triangular with these subquotients as diagonal blocks, proving (i)--(iii).

    Let $P_{j(s)}$, $Q_{j(s)}$ be the coordinate embedding and projection
    of the matched slot and set $V_s=G^{-1}P_{j(s)}X_s$,
    $W_s=X_s^{-1}Q_{j(s)}G$. Then $W_sV_t=X_s^{-1}Q_{j(s)}P_{j(t)}X_t$, and
    (\ref{eq:asym_block_biorthogonality}) gives (iv). For a nonempty word
    $w$, $B^w=G^{-1}T^wG$, and Lemma~\ref{lem:asym_triangular} gives
    $Q_{j(s)}T^wP_{j(s)}=X_sC_s^wX_s^{-1}$, so
    $W_sB^wV_s=X_s^{-1}Q_{j(s)}T^wP_{j(s)}X_s=C_s^w$, proving (v).

    For the remainder, $G\widehat C^iG^{-1}=\sum_{s\in S}P_{j(s)}X_sC_s^iX_s^{-1}Q_{j(s)}$
    is the block-diagonal matrix agreeing with $T^i$ on every matched
    diagonal block; every unmatched diagonal block of $T^i$ is already
    zero, so $GR^iG^{-1}=T^i-G\widehat C^iG^{-1}$ is strictly block upper
    triangular for the $r$-block decomposition. A nonzero entry of a
    product of $r$ strictly block upper-triangular $r$-block matrices would
    need a strictly increasing chain of $r+1$ block indices, which does not
    exist among $r$ positions; hence $GR^{i_1}\cdots R^{i_r}G^{-1}=0$,
    proving (vi). Finally, dimensions add along the composition series, and
    each zero-module step contributes dimension one, giving (vii).
\end{proof}

\begin{theorem}[Weighted canonical-form compression]
    \label{thm:asym_weighted_compression}
    \lean{MPSTensor.exists_weightedMultiBlockCompression_of_isNormal,
        MPSTensor.MultiBlockCompression.left_mul_evalWord_mul_right_smul}
    \leanok
    \uses{def:same_mpv2_pos, def:normal}
    Let $A=\bigoplus_{k=1}^g\bigoplus_{q=1}^{n_k}\mu_kA_{(k)}$ be a
    canonical target with each $A_{(k)}$ normal of bond dimension $D_k$
    and $\mu_k\in\C^\times$, and suppose
    $\ket{V^{(N)}(B)}=\ket{V^{(N)}(A)}$ for every $N>0$. For every block
    label $k$ and copy $q\in\{1,\ldots,n_k\}$, there are maps
    $V_{k,q}:\C^{D_k}\to\C^{D_B}$ and $W_{k,q}:\C^{D_B}\to\C^{D_k}$,
    mutually biorthogonal, with
    \begin{align}
        W_{k,q}B^{i_1}\cdots B^{i_N}V_{k,q}
         & =\mu_k^NA_{(k)}^{i_1}\cdots A_{(k)}^{i_N}
        \label{eq:asym_weighted_compression}
    \end{align}
    for every $N\geq1$ and every word.
\end{theorem}

\begin{proof}
    \leanok
    \uses{thm:asym_multiblock}
    Apply Theorem~\ref{thm:asym_multiblock} to the slots
    $C_{k,q}=\mu_kA_{(k)}$: for a word of length $N$,
    $C_{k,q}^{i_1}\cdots C_{k,q}^{i_N}=\mu_k^NA_{(k)}^{i_1}\cdots
        A_{(k)}^{i_N}$, so clause (v) of Theorem~\ref{thm:asym_multiblock}
    gives (\ref{eq:asym_weighted_compression}).
\end{proof}

% The Lean formalization of Theorem~\ref{thm:asym_multiblock} takes the
% hypothesis in the word-trace form $\tr(B^w)=\tr(A^w)$, and its conclusion
% chooses every matching isomorphism $X_s$ in clause (ii) to be the identity,
% which is stronger than the printed clause above; the underlying
% construction already produces this normalization because the composition
% series identifies each matched subquotient with the block's own
% word-algebra module, not merely with an isomorphic copy.

\begin{proposition}[Periodic data of a compression]
    \label{prop:asym_compression_trace}
    \lean{MPSTensor.MultiBlockCompression.trace_evalWord_eq_sum,
        MPSTensor.MultiBlockCompression.mpv_eq_sum}
    \leanok
    \uses{thm:asym_multiblock, lem:asym_triangular}
    Conversely, whenever a tensor $B$ admits the block triangular form of
    Theorem~\ref{thm:asym_multiblock}, clauses (i)--(iii), for a family
    $(C_s)_{s\in S}$, its periodic data are those of the direct sum:
    $\tr(B^w)=\sum_{s\in S}\tr(C_s^w)$ for every nonempty word $w$, that
    is, $\ket{V^{(N)}(B)}=\ket{V^{(N)}(\bigoplus_sC_s)}$ for every $N\geq1$.
\end{proposition}

\begin{proof}
    \leanok
    \uses{lem:asym_triangular}
    The trace is invariant under the change of basis and is the sum of the
    traces of the diagonal blocks; by Lemma~\ref{lem:asym_triangular} the
    diagonal blocks of $B^w$ are the words of the diagonal blocks, which are
    the $C_s^w$ on the matched slots and $0$ on the one-dimensional zero
    slots.
\end{proof}

\begin{proposition}[A vanishing remainder gives sitewise intertwiners]
    \label{prop:asym_split_intertwiners}
    \lean{MPSTensor.MultiBlockCompression.mul_right_eq_right_mul,
        MPSTensor.MultiBlockCompression.left_mul_eq_mul_left}
    \leanok
    \uses{thm:asym_multiblock}
    In the setting of Theorem~\ref{thm:asym_multiblock}, suppose that the
    remainder vanishes, $R^i=0$ for every physical index $i$. Then for
    every $i$ and every $s\in S$,
    \begin{align}
        B^iV_s & =V_sC_s^i, &
        W_sB^i & =C_s^iW_s.
    \end{align}
\end{proposition}

\begin{proof}
    \leanok
    Since $R^i=0$, $B^i=\sum_{t\in S}V_tC_t^iW_t$. Multiplying on the
    right by $V_s$, or on the left by $W_s$, and using the
    biorthogonality $W_tV_s=\delta_{ts}\id_{D_s}$ of clause (iv) leaves
    only the term $t=s$.
\end{proof}

\begin{proposition}[Projector-weighted direct sums]
    \label{prop:asym_projector_weighted_sum}
    \lean{MPSTensor.MultiBlockCompression.ofProjectorSum,
        MPSTensor.MultiBlockCompression.remainder_ofProjectorSum}
    \leanok
    \uses{thm:asym_multiblock, lem:asym_triangular}
    Let the bond space of $B$ be $\C^m\otimes\C^{D}$ and let $\Pi_1,\ldots$
    be mutually orthogonal rank-one idempotents of $\C^m$ cut out by an
    invertible matrix $U$ (the idempotents $U^{-1}E_{jj}U$). If every
    letter has the form
    $B^i=\sum_{s\in S}\mu_s\,\Pi_{j(s)}\otimes C_s^i$ with distinct
    positions $j(s)$, then $B$ admits the block triangular form of
    Theorem~\ref{thm:asym_multiblock} onto the weighted blocks $\mu_sC_s$,
    with $z=(m-|S|)D$ zero factors, the remainder $R^i=0$, and sitewise
    intertwiners $B^iV_s=V_s\,\mu_sC_s^i$ and $W_sB^i=\mu_sC_s^iW_s$ for
    every slot: the extension splits.
\end{proposition}

\begin{proof}
    \leanok
    \uses{lem:asym_triangular, prop:asym_split_intertwiners}
    Conjugating by $U\otimes\id$ turns each $\Pi_{j(s)}$ into a matrix
    unit, so every letter becomes block diagonal with the blocks
    $\mu_sC_s^i$ at the positions $j(s)$ and zero blocks elsewhere; block
    diagonality gives the triangular form and the vanishing remainder, and
    Proposition~\ref{prop:asym_split_intertwiners} gives the intertwining
    relations.
\end{proof}

%
\section{Products of matrix product operators}%
\label{sec:asym_mpo_product}

Let $M_\alpha$ and $M_\beta$ be matrix product operator tensors. Their
stacked product is the tensor
$B^{ij}=\sum_m M_\alpha^{im}\otimes M_\beta^{mj}$,
where the pair $(i,j)$ is treated as one physical index. Write $O_N(M)$
for the periodic operator of length $N$ generated by $M$.

A word in the paired indices specifies a matrix entry of $O_N(M)$.
That entry is the trace of the corresponding product of tensor matrices.
Thus identities between word traces at every positive length determine
identities between periodic operators.

\begin{lemma}[Word traces determine the periodic operator]%
    \label{lem:asym_mpo_trace_eq}
    \lean{MPOTensor.trace_evalWord_toMPSTensor,
        MPOTensor.mpo_eq_of_trace_evalWord_toMPSTensor,
        MPOTensor.mpo_eq_pow_smul_of_trace_evalWord_toMPSTensor,
        MPSTensor.trace_evalWord_smul}
    \leanok
    \uses{def:mpo_family, def:mpo_to_mps}
    The trace of a length-$N$ word in the paired indices of $M$ is the
    matrix entry of $O_N(M)$ at the ket and bra configurations specified
    by the word. Consequently, if two tensors $M$ and $\widetilde M$
    have equal word traces at every positive length, then
    $O_N(M)=O_N(\widetilde M)$ for every $N\geq1$. If instead, for a
    scalar $c$, every length-$N$ word trace of $M$ equals $c^N$ times
    the corresponding trace of $\widetilde M$, then
    $O_N(M)=c^NO_N(\widetilde M)$. Rescaling a tensor
    $A$ by $c$ multiplies each length-$N$ word trace by $c^N$.
\end{lemma}

\begin{proof}\leanok
    A word in paired indices specifies separate ket and bra words, and
    its trace is the corresponding entry of $O_N(M)$. Equality of these
    entries for all configurations gives equality of the operators.
    Rescaling a tensor by $c$ multiplies each letter's matrix by $c$,
    hence a length-$N$ word evaluation, and its trace, by $c^N$.
\end{proof}

\begin{theorem}[Reduction of a product of matrix product operators]
    \label{thm:asym_mpo_product}
    \lean{MPOTensor.exists_multiBlockCompression_mulTensor_of_mpo_eq_sum,
        MPOTensor.left_mul_evalWord_mul_right_smul}
    \leanok
    \uses{thm:asym_weighted_compression}
    Let $\{M_\gamma\}_\gamma$ be a finite family of matrix product
    operator tensors of positive bond dimension, each normal when its
    pair of physical indices is treated as one index. Suppose that
    nonzero scalars $\lambda_\gamma$ satisfy
    $O_N(B)=\sum_\gamma\lambda_\gamma^NO_N(M_\gamma)$ for every $N\geq1$.
    Then for every $\gamma$ there are maps $V_\gamma$ and $W_\gamma$ with
    $W_\gamma V_\gamma=\id$ and
    $W_\gamma B^{i_1j_1}\cdots B^{i_Nj_N}V_\gamma
        =\lambda_\gamma^NM_\gamma^{i_1j_1}\cdots M_\gamma^{i_Nj_N}$
    for every $N\geq1$ and every sequence of index pairs. The pairs may be
    chosen mutually biorthogonal, and the remainder after all matched
    sectors are removed is nilpotent.
\end{theorem}

\begin{center}
    % Formula: B^{ij}=\sum_mM_\alpha^{im}\otimes M_\beta^{mj}, the
    % stacked product defined in the paragraph preceding this theorem.
    % Source: \S\ref{sec:asym_mpo_product}, definition of the stacked
    % product.
    % Ink-to-index map: the upper row is $M_\alpha$ with physical legs
    % $i$, the lower row is $M_\beta$ with physical legs $j$; the shared
    % vertical pairleg at each site is the summed index $m$.
    % Contraction set: the vertical pairlegs (index $m$) at every site
    % and the horizontal virtual bonds within each row; the physical
    % legs $(i,j)$ remain open, read together as one letter of the pair
    % alphabet.
    % Boundary signature: (virtual-west=0 | virtual-east=0 |
    % physical-up=3 | physical-down=3) on both sides.
    \tenkzkernel
    \begin{tenkz}[rows={op,op}, cols=4, boundary=periodic]
        \tn[skin=mpo, ports={90:physical:$i_1$, 270:physical}]{M_\alpha} &
        \tn[skin=mpo, ports={90:physical:$i_2$, 270:physical}]{M_\alpha} &
        \tn[skin=dots, ports={180:virtual, 0:virtual}]{} &
        \tn[skin=mpo, ports={90:physical:$i_N$, 270:physical}]{M_\alpha} \\
        \tn[skin=mpo, ports={90:physical, 270:physical:$j_1$}]{M_\beta} &
        \tn[skin=mpo, ports={90:physical, 270:physical:$j_2$}]{M_\beta} &
        \tn[skin=dots, ports={180:virtual, 0:virtual}]{} &
        \tn[skin=mpo, ports={90:physical, 270:physical:$j_N$}]{M_\beta}
    \end{tenkz}
    \quad $=$ \quad
    \begin{tenkz}[rows={op}, cols=4, boundary=periodic]
        \tn[skin=mpo, ports={90:physical:$i_1$, 270:physical:$j_1$}]{B} &
        \tn[skin=mpo, ports={90:physical:$i_2$, 270:physical:$j_2$}]{B} &
        \tn[skin=dots, ports={180:virtual, 0:virtual}]{} &
        \tn[skin=mpo, ports={90:physical:$i_N$, 270:physical:$j_N$}]{B}
    \end{tenkz}
\end{center}

\begin{proof}
    \leanok
    \uses{thm:asym_weighted_compression}
    Treat each pair $(i,j)$ as one letter. The operator identity is the
    equality of the periodic vectors of $B$ and of $\bigoplus_\gamma
        \lambda_\gamma M_\gamma$ on the pair alphabet, so
    Theorem~\ref{thm:asym_weighted_compression} applies.
\end{proof}

\begin{corollary}[Fusion tensors of periodic matrix product operator
        representations]
    \label{cor:asym_fusion_tensors_group}
    \lean{MPOTensor.GroupFamily.IsNormalRepresentation.exists_fusionTensors}
    \leanok
    \uses{cor:asym_single_block, def:mpo_to_mps, def:normal,
        def:mpdo_operator_product}
    Let $\{U_g\}_{g\in G}$ be a matrix product operator representation of
    a group $G$, not necessarily unitary, by tensors $T_g$ of positive bond
    dimension whose doubled-index tensors are normal, with $U_gU_h=U_{gh}$
    on periodic chains of every positive length. Injective tensors are
    normal, so this covers the injective representations of
    \cite{GarreRubio2023MPOSym}; neither unitarity, simplicity, nor
    $U_e=\id$ is needed. Then for all $g,h\in G$
    there are fusion tensors $V$ and $W$ with $WV=\id$ and
    $W\,(T_gT_h)^{i_1j_1}\cdots(T_gT_h)^{i_Nj_N}\,V
        =T_{gh}^{i_1j_1}\cdots T_{gh}^{i_Nj_N}$
    for every $N\geq0$, and the remainder is nilpotent of order at most
    $D_gD_h-D_{gh}+1$.
\end{corollary}

\begin{proof}
    \leanok
    \uses{cor:asym_single_block}
    The stacked tensor of $T_g$ and $T_h$ generates the same periodic
    operators as $T_{gh}$, which is normal; apply
    Corollary~\ref{cor:asym_single_block}.
\end{proof}

\begin{corollary}[Action on matrix product vectors]
    \label{cor:asym_action_tensors}
    \lean{MPOTensor.exists_isReduction_actTensor_of_isNormal,
        MPOTensor.exists_multiBlockCompression_actTensor_of_isNormal}
    \leanok
    \uses{cor:asym_single_block, thm:asym_weighted_compression}
    Let $T$ be a matrix product operator tensor and $A_x$ a matrix product
    state tensor, and let $T\cdot A_x$ be the tensor of the operator applied
    to the state, $(T\cdot A_x)^i=\sum_jT^{ij}\otimes A_x^j$. If the periodic
    operators of $T$ map the periodic vectors of $A_x$ to those of a normal
    tensor $A_y$ of positive bond dimension at every positive length, then
    $T\cdot A_x$ reduces to $A_y$ with a nilpotent remainder. If, at every
    positive length, they map them to a finite weighted sum
    $\sum_s\mu_s^N\ket{V^{(N)}(A_s)}$ over normal tensors of positive bond
    dimension with nonzero weights $\mu_s$, then $T\cdot A_x$ admits a
    simultaneous reduction to the weighted blocks $\mu_sA_s$.
\end{corollary}

\begin{center}
    % Formula: (T\cdot A_x)^i=\sum_jT^{ij}\otimes A_x^j, the tensor
    % defined in the statement of this corollary, contracted into its
    % periodic state.
    % Source: Corollary~\ref{cor:asym_action_tensors}.
    % Ink-to-index map: the upper row is the operator tensor $T$ with
    % physical legs $i$, the lower row is the state tensor $A_x$; the
    % shared vertical pairleg at each site is the summed index $j$.
    % Contraction set: the vertical pairlegs (index $j$) at every site
    % and the horizontal virtual bonds within each row, both closed by
    % the periodic trace; only the upper physical legs of $T$ remain
    % open.
    % Boundary signature: (virtual-west=0 | virtual-east=0 |
    % physical-up=3 | physical-down=0).
    \tenkzkernel
    \begin{tenkz}[rows={op,ket}, cols=4, boundary=periodic]
        \tn[skin=mpo, ports={90:physical:$i_1$}]{T} &
        \tn[skin=mpo, ports={90:physical:$i_2$}]{T} &
        \tn[skin=dots, ports={180:virtual, 0:virtual}]{} &
        \tn[skin=mpo, ports={90:physical:$i_N$}]{T} \\
        \tn[label pos=s, ports={180:virtual, 0:virtual}]{A_x} &
        \tn[label pos=s, ports={180:virtual, 0:virtual}]{A_x} &
        \tn[skin=dots, ports={180:virtual, 0:virtual}]{} &
        \tn[label pos=s, ports={180:virtual, 0:virtual}]{A_x}
    \end{tenkz}
\end{center}

\begin{proof}
    \leanok
    \uses{cor:asym_single_block, thm:asym_weighted_compression}
    The periodic vectors of $T\cdot A_x$ are the images of those of $A_x$
    under the periodic operators of $T$, so the hypotheses of
    Corollary~\ref{cor:asym_single_block} and of
    Theorem~\ref{thm:asym_weighted_compression} hold.
\end{proof}

\begin{corollary}[Periodic identities from block reductions]
    \label{cor:asym_compression_periodic}
    \lean{MPOTensor.mpo_mul_eq_sum_of_multiBlockCompression,
        MPOTensor.mpo_mulVec_mpv_eq_sum_of_multiBlockCompression}
    \leanok
    \uses{prop:asym_compression_trace}
    If the stacked product of $M_\alpha$ and $M_\beta$ admits the block
    triangular form of Theorem~\ref{thm:asym_multiblock} for a family
    $\{C_s\}_{s\in S}$ of operator tensors, then
    $O_N(M_\alpha)O_N(M_\beta)=\sum_{s\in S}O_N(C_s)$ for every $N\geq1$.
    If the tensor $T\cdot A_x$ admits the same form for a family
    $\{C_s\}_{s\in S}$ of state tensors, then
    $O_N(T)\ket{V^{(N)}(A_x)}=\sum_{s\in S}\ket{V^{(N)}(C_s)}$ for every
    $N\geq1$.
\end{corollary}

\begin{proof}
    \leanok
    \uses{prop:asym_compression_trace}
    Apply Proposition~\ref{prop:asym_compression_trace} on the pair
    alphabet for the product and on the physical alphabet for
    $T\cdot A_x$. The definitions of the periodic operator and of
    $T\cdot A_x$ give the two identities.
\end{proof}

\subsection{Fusion and action tensors with multiplicity}%
\label{subsec:asym_fusion_multiplicity}

The fusion rules of a matrix product operator algebra carry
nonnegative integer multiplicities: $O_N(M_a)O_N(M_b)=\sum_c
    N_{ab}^cO_N(M_c)$ for every $N\geq1$ \cite{GarreRubio2023MPOSym}. Each copy
of a channel $c$ is then a separate block of the reduction.

\begin{theorem}[Reduction onto blocks with multiplicity]
    \label{thm:asym_multiplicity_reductions}
    \lean{MPSTensor.exists_multiBlockCompression_sigma_of_isNormal,
        MPSTensor.exists_multiplicityReductions_of_isNormal}
    \leanok
    \uses{thm:asym_multiblock}
    Let $\{C_c\}_c$ be a finite family of normal tensors of positive bond
    dimension and let $N_c\geq0$ be integers. If a tensor $B$ of bond
    dimension $D$ satisfies $\ket{V^{(N)}(B)}=\sum_cN_c\ket{V^{(N)}(C_c)}$
    for every $N\geq1$, then for every pair $(c,\mu)$ with $\mu<N_c$
    there are maps $V_{c,\mu}$ and $W_{c,\mu}$ such that
    $V_{c,\mu}W_{d,\nu}=\delta_{cd}\delta_{\mu\nu}\id$,
    $V_{c,\mu}B^{i_1}\cdots B^{i_N}W_{c,\mu}=C_c^{i_1}\cdots C_c^{i_N}$ for
    every $N\geq0$, and every product of at least $D$ matrices of the
    remainder $B^i-\sum_{c,\mu}W_{c,\mu}C_c^iV_{c,\mu}$ vanishes.
\end{theorem}

\begin{proof}
    \leanok
    \uses{thm:asym_multiblock}
    The word traces of $B$ are those of the family in which $C_c$ occurs
    $N_c$ times, so Theorem~\ref{thm:asym_multiblock} applies to that
    family. The number of blocks plus the number of zero slots is at most
    $D$, which bounds the nilpotency order.
\end{proof}

\begin{corollary}[Fusion and action tensors with multiplicity]
    \label{cor:asym_fusion_action_multiplicity}
    \lean{MPOTensor.exists_fusionTensors_of_mpo_mul_eq_sum,
        MPOTensor.IsMPOFusionAlgebra.exists_fusionTensors,
        MPOTensor.exists_actionTensors_of_mpo_mulVec_eq_sum}
    \leanok
    \uses{thm:asym_multiplicity_reductions}
    Let $\{M_c\}_c$ be matrix product operator tensors of positive bond
    dimensions $\chi_c$, each normal when its pair of physical indices is
    treated as one index, with
    $O_N(M_a)O_N(M_b)=\sum_cN_{ab}^cO_N(M_c)$ for every $N\geq1$. Then
    the stacked product of $M_a$ and $M_b$ has fusion tensors
    $V_{ab}^{c,\mu}$, $W_{ab}^{c,\mu}$, $\mu<N_{ab}^c$, as in
    Theorem~\ref{thm:asym_multiplicity_reductions}, with nilpotency order
    at most $\chi_a\chi_b$. Likewise, if
    $O_N(T)\ket{V^{(N)}(A_x)}=\sum_yM_y\ket{V^{(N)}(A_y)}$ for every
    $N\geq1$ with normal $A_y$ of positive bond dimension, then
    $T\cdot A_x$ has action tensors onto $M_y$ copies of each $A_y$.
\end{corollary}

\begin{proof}
    \leanok
    \uses{thm:asym_multiplicity_reductions}
    The periodic vectors of the stacked product on the pair alphabet are
    the matrix entries of $O_N(M_a)O_N(M_b)$, and those of $T\cdot A_x$
    are the entries of $O_N(T)\ket{V^{(N)}(A_x)}$; apply
    Theorem~\ref{thm:asym_multiplicity_reductions}.
\end{proof}

\begin{corollary}[Gauge freedom of a multiplicity-free fusion channel]
    \label{cor:asym_fusion_gauge_single}
    \lean{MPSTensor.IsReduction.exists_boundary_dressed_proportional_of_forall_le,
        MPOTensor.exists_boundary_dressed_proportional_of_mpo_mul_eq}
    \leanok
    \uses{thm:mps_reduction_boundary_dressed_uniqueness}
    Let $(V,W)$ and $(V',W')$ be two reductions of the stacked product $B$
    of $M_a$ and $M_b$ onto the same normal tensor $M_c$, each with a
    remainder all of whose products of at least $K$ factors vanish. Then
    there is $z\neq0$ with $VB^{w}=z\,V'B^{w}$ and
    $B^{w}W=z^{-1}B^{w}W'$ for every word $w$ of length greater than
    $2K$. When $O_N(M_a)O_N(M_b)=O_N(M_c)$ for every $N\geq1$,
    Corollary~\ref{cor:asym_fusion_action_multiplicity} provides such
    reductions with $K=\chi_a\chi_b$.
\end{corollary}

\begin{proof}
    \leanok
    \uses{thm:mps_reduction_boundary_dressed_uniqueness}
    Apply Theorem~\ref{thm:mps_reduction_boundary_dressed_uniqueness} on
    the pair alphabet.
\end{proof}
%
\section{Structure coefficients at every positive length}
\label{sec:asym_unblocked_coefficients}

The matrix product operator algebras of renormalization fixed points may have
structure coefficients that depend on the chain length. The result below gives
a general power-sum constraint: the coefficients in an expansion over
inequivalent normal tensors are power sums of fixed finite multisets of complex
weights, at every positive length and with no blocking.
Normality means injectivity after some blocking (Definition~\ref{def:normal}),
and separation excludes gauge equivalence up to any nonzero complex scalar.

\begin{theorem}[Uniqueness of nonzero weights from tail power sums]
    \label{thm:asym_tail_weights}
    \lean{Multiset.eq_of_notMem_zero_of_forall_sum_map_pow_eq,
        Multiset.eq_zero_of_notMem_zero_of_forall_sum_map_pow_eq_zero}
    \leanok
    Let $\Lambda$ and $\Lambda'$ be finite multisets in $\C\setminus\{0\}$.
    If their power sums agree for all sufficiently large integers $N$, then
    $\Lambda=\Lambda'$, including multiplicities. In particular, if the
    power sums of $\Lambda$ vanish for all sufficiently large $N$, then
    $\Lambda$ is empty.
\end{theorem}
\begin{proof}
    \leanok
    List the distinct entries of the two multisets as
    $z_0,\ldots,z_{r-1}$, and let $b_j$ be the difference of their
    multiplicities. For a threshold $N_0$ and each $0\leq i<r$,
    \begin{align}
        \sum_{j=0}^{r-1}(b_jz_j^{N_0})z_j^i & = 0.
    \end{align}
    The Vandermonde matrix is invertible, so $b_jz_j^{N_0}=0$.
    Since $z_j\neq0$, all multiplicities agree. Comparison with the empty
    multiset proves the vanishing assertion.
\end{proof}

\begin{theorem}[Power-sum coefficients without blocking]
    \label{thm:asym_unblocked_coefficients}
    \lean{MPSTensor.exists_unblocked_powerSum_coeff_unique,
        MPSTensor.exists_unblocked_powerSum_coeff}
    \leanok
    \uses{def:normal}
    Let $B$ be an MPS tensor of bond dimension $D_B$, and let
    $\{M_\gamma\}_{\gamma\in\Gamma}$ be a finite family of normal MPS
    tensors on the same physical alphabet, with $D_\gamma>0$.
    For distinct $\gamma,\delta$ of equal bond dimension, assume there
    are no $X\in\GL(D_\gamma,\C)$ and $z\in\C\setminus\{0\}$ such that
    $M_\delta^i=zXM_\gamma^iX^{-1}$ for every physical index $i$.
    Suppose that for every $N>0$ there are coefficients $c_\gamma^{(N)}$
    satisfying
    \begin{align}
        \ket{V^{(N)}(B)}
         & =\sum_{\gamma\in\Gamma}c_\gamma^{(N)}
        \ket{V^{(N)}(M_\gamma)}.
        \label{eq:asym_unblocked_span}
    \end{align}
    Then there are integers $n_\gamma\geq0$ and nonzero complex weights
    $\mu_{\gamma k}$, indexed by $0\leq k<n_\gamma$, such that for every
    $N>0$,
    \begin{align}
        \ket{V^{(N)}(B)}
         & =\sum_{\gamma\in\Gamma}
        \left(\sum_{k=0}^{n_\gamma-1}\mu_{\gamma k}^{N}\right)
        \ket{V^{(N)}(M_\gamma)}.
        \label{eq:asym_unblocked_expansion}
    \end{align}
    The multiplicities satisfy $\sum_{\gamma\in\Gamma}n_\gamma D_\gamma\leq D_B$.
    There is a positive threshold $N_0$ such that, for every $N\geq N_0$,
    any coefficients satisfying~(\ref{eq:asym_unblocked_span}) obey
    $c_\gamma^{(N)}=\sum_{k=0}^{n_\gamma-1}\mu_{\gamma k}^{N}$.
    The weight multisets are unique: any finite families of nonzero
    weights giving an expansion of the form~(\ref{eq:asym_unblocked_expansion})
    at all sufficiently large lengths have, for each $\gamma$, exactly
    the same entries with the same multiplicities. Empty weight
    multisets are allowed.
\end{theorem}
\begin{proof}
    \leanok
    \uses{thm:irred_decomp,thm:periodic_overlap_dichotomy,
        thm:periodic_basis_eventually_li,thm:asym_tail_weights}
    Apply Theorem~\ref{thm:irred_decomp}, normalize each nonzero irreducible
    factor to a periodic tensor, and group factors whose periodic vectors
    are proportional by a fixed nonzero scalar raised to the length.
    The irreducible-form framework is described in
    Ref.~\cite{DeLasCuevas2017Irreducible}. This gives separated periodic
    representatives $A_j$ of dimensions $d_j$, positive multiplicities
    $r_j$, and nonzero weights $a_{jq}$ such that, for every $N>0$,
    \begin{align}
        \ket{V^{(N)}(B)}
                      & =\sum_j\left(\sum_{q=0}^{r_j-1}a_{jq}^{N}\right)
        \ket{V^{(N)}(A_j)}, \label{eq:asym_periodic_expansion}           \\
        \sum_j r_jd_j & \leq D_B.
    \end{align}
    Grouping preserves the dimension count. Indeed, if periodic tensors
    $A,A'$ of periods $m,m'$ have such proportional vectors, then
    $|O_{AA'}(N)|^2=O_{AA}(N)O_{A'A'}(N)$.
    Different dimensions would force the left side to tend to zero by
    Theorem~\ref{thm:periodic_overlap_dichotomy}, whereas along common
    multiples of the periods the right side tends to $mm'>0$.

    Normalize the targets to period-one tensors $C_\gamma$ and combine
    them with the representatives not related to any target by an invertible gauge
    and a nonzero scalar.
    The combined family has no repeated blocks. On sufficiently large
    multiples $pn$ of a common period, it is linearly independent by
    Theorem~\ref{thm:periodic_basis_eventually_li}.
    Subtracting~(\ref{eq:asym_unblocked_span}) from
    (\ref{eq:asym_periodic_expansion}), after absorbing matched terms
    into the targets, forces every unmatched representative to satisfy
    $\sum_{q=0}^{r_j-1}(a_{jq}^{p})^n=0$ for all sufficiently large $n$.
    Theorem~\ref{thm:asym_tail_weights} contradicts $r_j>0$.
    Thus every representative matches a target.

    Choose its matching label $\gamma(j)$ and a nonzero $\theta_j$ so that
    $\ket{V^{(N)}(M_{\gamma(j)})}
        =\theta_j^N\ket{V^{(N)}(A_j)}$.
    For each target, collect the weights $a_{jq}/\theta_j$ over its
    matching representatives. This proves~(\ref{eq:asym_unblocked_expansion}).
    Since $d_j=D_{\gamma(j)}$, the dimension bound follows.
    Normalize the normal targets and use their eventual linear independence
    to compare arbitrary expansion coefficients with these power sums.
    The nonzero normalization scalars can be cancelled, giving a positive
    threshold for coefficient uniqueness. Theorem~\ref{thm:asym_tail_weights}, applied to each
    target, then proves uniqueness against any tail expansion.
\end{proof}

\begin{theorem}[Power-sum structure coefficients for operator closure]
    \label{thm:asym_operator_coefficients}
    \lean{MPOTensor.exists_operatorProduct_powerSum_coeff_unique}
    \leanok
    \uses{def:normal}
    Let $\{M_\gamma\}_{\gamma\in\Gamma}$ be a finite family of MPO
    tensors of common physical dimension $d$, with $D_\gamma>0$. Regard the pair of physical indices
    $(i,j)$ as one MPS index. Assume these MPS tensors are normal and,
    for distinct labels of equal bond dimension, are not related by
    $M_\delta^{ij}=zXM_\gamma^{ij}X^{-1}$ for all $i,j$, with $z\neq0$
    and $X$ invertible. Fix $\alpha,\beta\in\Gamma$ and write $O_L(M)$
    for the periodic operator generated by $M$. Suppose
    $O_L(M_\alpha)O_L(M_\beta)$ belongs to
    $\spn_{\C}\{O_L(M_\gamma):\gamma\in\Gamma\}$ for every $L>0$.

    There are complex diagonal matrices $\chi_{\alpha,\beta,\gamma}$
    of sizes $n_\gamma\geq0$, with nonzero diagonal entries, such that
    for every $L>0$,
    \begin{align}
        O_L(M_\alpha)O_L(M_\beta)
         & =\sum_{\gamma\in\Gamma}
        \tr(\chi_{\alpha,\beta,\gamma}^{L})O_L(M_\gamma).
        \label{eq:asym_operator_expansion}
    \end{align}
    Their sizes satisfy $\sum_{\gamma\in\Gamma}n_\gamma D_\gamma
        \leq D_\alpha D_\beta$.
    Beyond a positive threshold, any coefficients of this operator
    expansion equal $c_{\alpha,\beta,\gamma}^{(L)}
        =\tr(\chi_{\alpha,\beta,\gamma}^{L})$.
    Any other finite families of nonzero diagonal entries giving such an
    expansion at all sufficiently large lengths agree target by target,
    including multiplicities.

    This uses the structure-coefficient notation of
    Ref.~\cite[Section~4.5]{Cirac2016MPDO_arXiv}, but asserts neither
    positive weights nor the idempotent condition from that source.
    No renormalization fixed-point channels or anomaly conclusion are inferred.
\end{theorem}
\begin{proof}
    \leanok
    \uses{thm:mpdo_operator_product_mpo,thm:asym_unblocked_coefficients}
    Form the product tensor $P^{ij}=\sum_s M_\alpha^{is}\otimes M_\beta^{sj}$,
    of bond dimension $D_\alpha D_\beta$.
    Theorem~\ref{thm:mpdo_operator_product_mpo} gives
    $O_L(P)=O_L(M_\alpha)O_L(M_\beta)$.
    Equality of operators is equivalent, entry by entry, to equality of
    MPS coefficients on the combined physical alphabet.
    Apply Theorem~\ref{thm:asym_unblocked_coefficients} to $P$ and take
    $\chi_{\alpha,\beta,\gamma}$ to have diagonal entries
    $\{\mu_{\gamma k}\}_{k=0}^{n_\gamma-1}$.
    The same entrywise equivalence transfers the expansion, eventual
    coefficient uniqueness, dimension bound, and uniqueness of nonzero
    entries to operators.
\end{proof}
%
\section{The splitting criterion}%
\label{sec:asym_splitting}

Periodic data alone does not in general imply the existence of sitewise
fusion tensors (Remark~\ref{rem:asym_zipper}). The following criterion gives
a sufficient algebraic condition for the compressions of
Theorem~\ref{thm:asym_multiblock} to split into sitewise intertwiners.

\begin{theorem}[Splitting criterion]
    \label{thm:asym_splitting}
    \lean{MPSTensor.exists_multiBlockCompression_remainder_eq_zero_of_star,
        MPSTensor.exists_multiBlockCompression_remainder_eq_zero_of_isSemisimpleModule,
        MPSTensor.isSemisimpleModule_wordModule_of_conjTranspose_mem_adjoin}
    \leanok
    \uses{thm:asym_multiblock, def:asym_word_module}
    In the setting of Theorem~\ref{thm:asym_multiblock}, suppose that the
    unital algebra generated by the matrices $B^i$ is closed under conjugate
    transposition, that is, $(B^i)^\dagger$ lies in the unital algebra generated
    by the $B^j$ for every $i$. Then the word-algebra module of $B$ is
    semisimple, and $B$ admits a block triangular form as in
    Theorem~\ref{thm:asym_multiblock} whose remainder vanishes: the block
    triangular form is block diagonal.
\end{theorem}

\begin{proof}
    \leanok
    \uses{thm:asym_multiblock, thm:asym_flag_theorem}
    Since the algebra is closed under conjugate transposition, the
    orthogonal complement of an invariant subspace is invariant. Thus
    every invariant subspace has an invariant complement and the module
    is semisimple. Every step of the flag of
    Theorem~\ref{thm:asym_flag_theorem} then has an invariant complement,
    the module is the direct sum of these complements, and a basis
    adapted to the direct sum makes every $B^i$ block diagonal with the
    blocks prescribed by the flag.
\end{proof}

\begin{corollary}[Sitewise fusion tensors]
    \label{cor:asym_fusion_tensors_star}
    \lean{MPSTensor.exists_fusionTensors_of_star}
    \leanok
    \uses{thm:asym_splitting}
    Under the hypotheses of Theorem~\ref{thm:asym_splitting} there are
    maps $V_s$ and $W_s$ with $W_sV_t=\delta_{st}\id$, $W_sB^wV_s=C_s^w$
    for every word $w$, and the sitewise relations $B^iV_s=V_sC_s^i$ and
    $W_sB^i=C_s^iW_s$ for every physical index $i$ and block occurrence $s$.
\end{corollary}

\begin{center}
    % Formula: B^iV_s=V_sC_s^i and W_sB^i=C_s^iW_s for every letter i.
    % Source: Corollary~\ref{cor:asym_fusion_tensors_star}.
    % Ink-to-index map: the ring is the compression map $V_s$ or $W_s$;
    % the square node is $B$ or its target $C_s$.
    % Contraction set: the single virtual bond joining the ring to the
    % square node in each pair; the physical leg $i$ and the two outer
    % virtual legs remain open.
    % Boundary signature: (virtual-west=1 | virtual-east=1 |
    % physical-up=1 | physical-down=0) for each pictured equality.
    \tenkzkernel
    \begin{tenkz}[rows={wire}, cols=2, west=open, east=open]
        \tn[label pos=s, ports={90:physical:$i$}]{B} & \tn[skin=ring]{V_s}
    \end{tenkz}
    \quad $=$ \quad
    \begin{tenkz}[rows={wire}, cols=2, west=open, east=open]
        \tn[skin=ring]{V_s} & \tn[label pos=s, ports={90:physical:$i$}]{C_s}
    \end{tenkz}
    \\[2ex]
    \begin{tenkz}[rows={wire}, cols=2, west=open, east=open]
        \tn[skin=ring]{W_s} & \tn[label pos=s, ports={90:physical:$i$}]{B}
    \end{tenkz}
    \quad $=$ \quad
    \begin{tenkz}[rows={wire}, cols=2, west=open, east=open]
        \tn[label pos=s, ports={90:physical:$i$}]{C_s} & \tn[skin=ring]{W_s}
    \end{tenkz}
\end{center}

\begin{proof}
    \leanok
    \uses{thm:asym_splitting}
    With vanishing remainder, $B^i=\sum_sV_sC_s^iW_s$, and the sitewise
    relations follow from the biorthogonality of the pairs.
\end{proof}

\begin{corollary}[Obstruction to closure under conjugate transposition]
    \label{cor:asym_anomaly_not_star}
    \lean{MPSTensor.not_forall_conjTranspose_mem_adjoin_of_forall_right_intertwiner_eq_zero,
        CZXCompression.czxSquare_not_starClosed}
    \leanok
    \uses{thm:asym_splitting}
    If, in the setting of Theorem~\ref{thm:asym_multiblock}, some block
    occurrence $s$ admits no nonzero sitewise right intertwiner $X$ with
    $B^iX=XC_s^i$ for all $i$, then the unital algebra generated by the
    $B^i$ is not closed under conjugate transposition. In particular, the
    stacked product of the anomalous CZX symmetry with itself, for which
    both sitewise intertwiner spaces vanish, generates a unital algebra
    that is not closed under conjugate transposition.
\end{corollary}

\begin{proof}
    \leanok
    \uses{thm:asym_splitting}
    By Theorem~\ref{thm:asym_splitting}, closure under conjugate
    transposition would give a compression pair whose right map is a
    nonzero sitewise intertwiner.
\end{proof}

\begin{remark}
    \label{rem:asym_anomaly_unitary}
    The physical operator of the CZX symmetry, a product of controlled-Z
    gates and single-site spin flips, is unitary at every length. The
    failure of closure under conjugate transposition in
    Corollary~\ref{cor:asym_anomaly_not_star} therefore concerns the
    virtual algebra of the stacked tensor only, not the physical operator;
    unitarity of the operator is not part of the formal statement.
    This algebraic conclusion does not compute the anomaly class.
    Determining that class requires a separate computation of the
    $3$-cocycle defined by the associator of the fusion tensors.
\end{remark}

\section{The single-block reduction of Moln\'ar--Ge--Schuch--Cirac}%
\label{sec:asym_single_block}

\begin{corollary}[Single-block case: the normal-target reduction]
    \label{cor:asym_single_block}
    \lean{MPSTensor.exists_isReduction_and_nilpotent_of_isNormal}
    \leanok
    \uses{def:same_mpv2_pos, def:normal, def:mps_rectangular_reduction}
    Let $A$ be a normal tensor (NT) of bond dimension $D_A\geq1$, let $B$
    be an arbitrary tensor of bond dimension $D_B$, and suppose that
    $\ket{V^{(N)}(B)}=\ket{V^{(N)}(A)}$ for every $N\geq1$.
    Then there are maps $V:\C^{D_A}\to\C^{D_B}$ and
    $W:\C^{D_B}\to\C^{D_A}$ with $WV=\id$ such that, for every $N\geq0$
    and every word of length $N$,
    \begin{align}
        WB^{i_1}\cdots B^{i_N}V & =A^{i_1}\cdots A^{i_N},
        \label{eq:asym_single_block_compression}
    \end{align}
    and the remainder $R^i=B^i-VA^iW$ generates a nilpotent non-unital
    algebra. Explicitly, for every word of length $r=D_B-D_A+1$,
    \begin{align}
        R^{i_1}R^{i_2}\cdots R^{i_r} & =0.
        \label{eq:asym_single_block_nilpotency}
    \end{align}
    In particular $D_B\geq D_A$.
\end{corollary}

\begin{center}
    % Formula: WB^{i_1}\cdots B^{i_N}V=A^{i_1}\cdots A^{i_N}
    % (\ref{eq:asym_single_block_compression}).
    % Source: Corollary~\ref{cor:asym_single_block}.
    % Ink-to-index map: the two rings are the single compression maps
    % $W$ (left) and $V$ (right), the special case of $W_s$, $V_s$ at
    % the unique slot; the interior nodes are copies of $B$.
    % Contraction set: every horizontal virtual bond between consecutive
    % copies of $B$ and between each $B$ and its adjacent ring; the
    % physical legs and the two outer virtual legs (the
    % $D_A$-dimensional indices of the matrix $A^{i_1}\cdots A^{i_N}$)
    % remain open.
    % Boundary signature: (virtual-west=1 | virtual-east=1 |
    % physical-up=2 | physical-down=0) for the finite schematic shown,
    % with the two explicit lettered nodes carrying the drawn physical
    % ports.
    \tenkzkernel
    \begin{tenkz}[rows={wire}, cols=5, west=open, east=open]
        \tn[skin=ring]{W} &
        \tn[label pos=s, ports={90:physical:$i_1$}]{B} &
        \tn[skin=dots]{} &
        \tn[label pos=s, ports={90:physical:$i_N$}]{B} &
        \tn[skin=ring]{V}
    \end{tenkz}
    \quad $=$ \quad
    \begin{tenkz}[rows={wire}, cols=3, west=open, east=open]
        \tn[label pos=s, ports={90:physical:$i_1$}]{A} &
        \tn[skin=dots]{} &
        \tn[label pos=s, ports={90:physical:$i_N$}]{A}
    \end{tenkz}
\end{center}

\begin{proof}
    \leanok
    \uses{thm:asym_multiblock}
    Apply Theorem~\ref{thm:asym_multiblock} to the single-block target
    $C=A$. Clauses (iv) and (v) give the compression pair and
    (\ref{eq:asym_single_block_compression}), clause (vii) gives
    $D_B=D_A+z\geq D_A$, and clause (vi) gives
    (\ref{eq:asym_single_block_nilpotency}) with $r=|S|+z=1+(D_B-D_A)$.
\end{proof}

The proportional case, in which $\ket{V^{(N)}(B)}=\lambda^N\ket{V^{(N)}(A)}$
for a scalar $\lambda\in\C^\times$, reduces to
Corollary~\ref{cor:asym_single_block} by replacing $B$ with
$\lambda^{-1}B$, and then $WB^{i_1}\cdots B^{i_N}V=\lambda^NA^{i_1}\cdots A^{i_N}$
with remainder $R^i=B^i-\lambda\,VA^iW$; the compression pair of that case
is recorded in Theorem~\ref{thm:mps_proportional_reduction_exists}.

Corollary~\ref{cor:asym_single_block} together with the proportional case
corresponds to Propositions~20 and~21 of~\cite{Molnar2018SemiInjective} (with
$V,W$ swapped), with three source corrections and an explicit nilpotency bound.
First, the hypothesis $\lambda\ne0$ is required because the printed
Proposition~20 fails at $\lambda=0$: the zero tensor $B=0$ satisfies the
hypothesis with no compression onto a nonzero normal target. Second, the
source proof silently normalizes $\lambda=1$; the reduction to the equality case
makes the scalar explicit. Both corrections are recorded in
\cite{gap:mgsc18_reduction_proportionality_scalar}
(\texttt{docs/paper-gaps/mgsc18_reduction_proportionality_scalar.tex}). Third, the remainder
algebra is generated without unit (a unital algebra contains a nonzero identity
and is not nilpotent). The triangular gauge of Theorem~\ref{thm:asym_multiblock}
yields the explicit nilpotency length bound $r=D_B-D_A+1$; see
\cite{gap:mgsc18_nilpotency_length_one_terminology}
(\texttt{docs/paper-gaps/mgsc18_nilpotency_length_one_terminology.tex}). The uniqueness statement
(Theorem~22 of~\cite{Molnar2018SemiInjective}) is not addressed here and fails
verbatim in the multi-block setting, as the matching isomorphisms $X_s$ of
Theorem~\ref{thm:asym_multiblock} are not unique in general.
%
\section{Rectangular reductions and residual algebras}%
\label{sec:asym_reductions}

This section develops the single-block rectangular reductions and residual
algebras of \cite[Propositions~20 and~21]{Molnar2018SemiInjective},
connecting positive-length periodic vector equality to compression pairs and
bounding the nilpotency length of the residual algebra.

\begin{definition}[Coefficient-dual inverse and reduction cross matrix]
    \label{def:mps_reduction_cross_matrix}
    \lean{MPSTensor.coefficientDualInverse,
        MPSTensor.reductionCrossMatrix}
    \leanok
    \uses{def:mps_tensor, def:injective}
    Let $\widetilde A=(\widetilde A^i)_i$ be injective, with bond dimension
    $D_A$, and choose the coefficient dual $C=(C^i)_i$ from a right inverse
    to the linear-combination map of $\widetilde A$.  Thus
    $C^i_{a'a}$ is the coefficient of $\widetilde A^i$ in the chosen
    expansion of the matrix unit $\lvert a\rangle\!\langle a'\rvert$.
    For a tensor $\widetilde B$ of bond dimension $D_B$, define the
    mixed-bond matrix on $\C^{D_B}\otimes\C^{D_A}$ by
    \begin{align}
        K(\widetilde B,C)
         & =\sum_i \widetilde B^i\otimes(C^i)^T.
        \notag
    \end{align}
    The superscript $T$ is the ordinary transpose, not the conjugate
    transpose.  This is the matrix drawn in
    \cite[proof of Proposition~20, lines~3817--3838]{Molnar2018SemiInjective}.
\end{definition}

\begin{lemma}[Coefficient dual and reversed-word power formulas]
    \label{lem:mps_reduction_cross_matrix_power}
    \lean{MPSTensor.coefficientDualInverse_sum_entry,
        MPSTensor.reductionCrossMatrix_apply,
        MPSTensor.reductionCrossMatrix_pow_eq_sum,
        MPSTensor.reductionCrossMatrix_pow_apply,
        MPSTensor.reductionCrossMatrix_coefficientDualInverse}
    \leanok
    \uses{def:mps_reduction_cross_matrix}
    The coefficient dual obeys the exact matrix-unit identity
    \begin{align}
        \sum_i \widetilde A^i_{xy}C^i_{a'a}
         & =\delta_{x,a}\delta_{y,a'}.
        \notag
    \end{align}
    Consequently,
    \begin{align}
        (K(\widetilde B,C)^n)_{(b,a),(b',a')}
         & =\sum_{i_1,\ldots,i_n}
        (\widetilde B^{i_1}\cdots\widetilde B^{i_n})_{bb'}
        (C^{i_n}\cdots C^{i_1})_{a'a}.
        \notag
    \end{align}
    In particular, pairing $\widetilde A$ with its own dual gives the
    rank-one matrix whose entry is
    $\delta_{x,a}\delta_{y,a'}$.  The reversed order
    $C^{i_n}\cdots C^{i_1}$ is forced by the transpose and records the
    upper-leg orientation in \cite[lines~3817--3865]{Molnar2018SemiInjective}.
\end{lemma}

\begin{proof}
    \leanok
    \uses{def:mps_reduction_cross_matrix}
    The coefficient identity follows by applying the chosen right inverse to a
    matrix unit.  Expanding the power of the finite sum defining $K$,
    multiplicativity of the Kronecker product collects the lower factors in
    site order.  Transposing the upper product reverses its factors, which gives
    the displayed entry formula.  Specializing the lower tensor to
    $\widetilde A$ and applying the coefficient identity gives the stated
    rank-one matrix.
\end{proof}

\begin{theorem}[Cross trace powers and source open-boundary contraction]
    \label{thm:mps_reduction_cross_trace_open}
    \lean{MPSTensor.SameMPV₂Pos.trace_evalWord,
        MPSTensor.trace_reductionCrossMatrix_pow_of_sameMPV₂Pos,
        MPSTensor.reductionOpenBoundaryMatrixContraction,
        MPSTensor.reductionOpenBoundaryContraction,
        MPSTensor.reductionOpenBoundaryMatrixContraction_eq_cross_sum,
        MPSTensor.reductionOpenBoundaryMatrixContraction_self,
        MPSTensor.reductionOpenBoundaryContraction_self,
        MPSTensor.reductionOpenBoundaryContraction_of_sameMPV₂Pos,
        MPSTensor.reductionOpenBoundaryMatrixContraction_blockTensor_of_sameMPV₂Pos}
    \leanok
    \uses{def:same_mpv2_pos, def:injective,
        def:mps_reduction_cross_matrix, thm:same_mpv2_pos_block_tensor}
    Suppose $\widetilde A$ and $\widetilde B$ have equal positive-length
    MPV coefficients and $\widetilde A$ is injective.  This equality hypothesis
    is the $\lambda=1$ specialization of the proportional premise printed in
    Proposition~20; the unrestricted unscaled conclusion is false
    \cite{gap:mgsc18_reduction_proportionality_scalar}.  Equality transports
    the trace of every nonempty word.  For every $n>0$,
    \begin{align}
        \tr K(\widetilde B,C)^n
         & =D_A^n.
        \notag
    \end{align}

    For arbitrary tensors $B_0,C_0$, every $n\geq0$, and every boundary matrix
    $X$, define
    \begin{align}
        \mathcal O_n(B_0,C_0;X)_{a'a}
         & =\sum_{i_1,\ldots,i_n}
        \tr\!\left(B_0^{i_1}\cdots B_0^{i_n}X\right)
        (C_0^{i_n}\cdots C_0^{i_1})_{a'a}.
        \notag
    \end{align}
    If $B_0$ has bond dimension $D_B$ and $C_0$ has bond dimension $D_A$,
    its exact cross-sum expansion is
    \begin{align}
        \mathcal O_n(B_0,C_0;X)_{a'a}
         & =\sum_{x,y}K(B_0,C_0)^n_{(x,a),(y,a')}X_{yx}.
        \notag
    \end{align}
    For $n>0$, if $B_0=\widetilde A$, $C_0=C$, and
    $X\in\C^{D_A\times D_A}$, then
    \begin{align}
        \mathcal O_n(\widetilde A,C;X)_{a'a}
         & =D_A^{n-1}X_{a'a}.
        \notag
    \end{align}
    The same-alphabet open-boundary contraction is the specialization
    $X=B_0^w$.  Thus $X=\widetilde A^w$ gives the target self-evaluation, while
    equality of the positive-length MPV coefficients gives
    \begin{align}
        \mathcal O_n(\widetilde B,C;\widetilde B^w)_{a'a}
         & =D_A^{n-1}(\widetilde A^w)_{a'a}.
        \notag
    \end{align}

    For the padded blocked contraction, suppose additionally that the original
    tensors $A$ and $B$ have equal positive-length MPV coefficients.  If
    $\widetilde A$ and $\widetilde B$ are obtained by blocking $p>0$ original
    sites, then for every unblocked lower tail word $w$,
    \begin{align}
        \mathcal O_n(\widetilde B,C;B^w)_{a'a}
         & =D_A^{n-1}(A^w)_{a'a}.
        \notag
    \end{align}
    The trace-power identity is the closed contraction at
    \cite[proof of Proposition~20, lines~3839--3865]{Molnar2018SemiInjective},
    and the open-boundary identity is the padded contraction at lines 3866--3887.
    The source next uses a rank-one rewrite and comparison to derive
    $VB^{i_1}\cdots B^{i_m}W=A^{i_1}\cdots A^{i_m}$ at lines 3888--3938.
\end{theorem}

\begin{proof}
    \leanok
    \uses{lem:mps_reduction_cross_matrix_power, lem:eval_word_block}
    Expand $K^n$ by Lemma~\ref{lem:mps_reduction_cross_matrix_power} and use
    the trace formula for a Kronecker product.  Positive-length MPV equality
    replaces the trace of each nonempty $\widetilde B$ word by the
    corresponding $\widetilde A$ trace.  The resulting scalar sum is therefore
    the trace of the target cross matrix power.  The coefficient-dual identity
    gives that target cross matrix its rank-one form, whose $n$th trace is
    $D_A^n$.

    Expanding the lower trace into matrix entries gives the cross-sum formula
    for arbitrary $X$.  For the target tensor, the rank-one form of
    $K(\widetilde A,C)$ evaluates this sum as $D_A^{n-1}X_{a'a}$.
    Specializing $X$ to a word evaluation gives the self-evaluation formula.

    For blocked tensors, Lemma~\ref{lem:eval_word_block} identifies each
    blocked word with the evaluation of its flattened word.  Since $n,p>0$,
    this flattened word is nonempty.  Append the unblocked tail $w$, apply
    positive-length MPV equality to the resulting original-site word, and use
    the target self-evaluation with $X=A^w$.
\end{proof}

\begin{theorem}[Rank-one factorizations]
    \label{thm:rank_one_factorizations}
    \lean{LinearMap.exists_smulRight_of_finrank_range_eq_one,
        Matrix.exists_eq_vecMulVec_of_rank_eq_one,
        Matrix.exists_rectangular_factors_of_rank_eq_one}
    \leanok
    Let $f\colon X\to Y$ be a linear map over a division ring $K$ whose range
    has dimension one.  Then there are a nonzero linear functional
    $\varphi\colon X\to K$, a nonzero vector $y\in Y$, and a vector $x\in X$
    such that
    \begin{align}
        \varphi(x) & =1, & f(x) & =y, & f(z) & =\varphi(z)y
        \quad\text{for every }z\in X.
        \notag
    \end{align}
    Consequently, if $I$ and $J$ are finite index sets, then every rank-one
    matrix $M\in K^{I\times J}$ over a field is an outer product of two
    nonzero vectors: there are $w\in K^I$ and $v\in K^J$ such that
    \begin{align}
        M_{ij} & =w_i v_j.
        \notag
    \end{align}
    In particular, if both indices of $M$ are the product
    $D_B\times D_A$, then there are nonzero rectangular matrices
    $W\in K^{D_B\times D_A}$ and $V\in K^{D_A\times D_B}$ such that
    \begin{align}
        M_{(b,a),(b',a')}
         & =W_{ba}V_{a'b'}.
        \notag
    \end{align}
    This is the rank-one factorization used in
    \cite[proof of Proposition~20, lines~3866--3936]{Molnar2018SemiInjective}.
\end{theorem}

\begin{proof}
    \leanok
    Choose an isomorphism from the scalar division ring to the one-dimensional
    range of $f$.  Composing its inverse with $f$ gives $\varphi$, while the
    image of $1$ gives $y$.  Choose a preimage $x$ of $y$; the construction gives
    $\varphi(x)=1$, $f(x)=y$, and the stated factorization of $f$.

    In standard bases this factorization is an outer product of two nonzero
    coordinate vectors.  For the product-index specialization, reshape the
    left vector as a $D_B\times D_A$ matrix and the right vector as a
    $D_A\times D_B$ matrix, preserving the displayed index order.
\end{proof}


\begin{theorem}[Equality-case existence of rectangular reductions]
    \label{thm:mps_equality_reduction_exists}
    \lean{MPSTensor.exists_isReduction_of_isInjective_of_sameMPV₂Pos,
        MPSTensor.exists_isReduction_of_isNormal_of_sameMPV₂Pos}
    \leanok
    \uses{def:mps_rectangular_reduction, def:same_mpv2_pos, def:injective,
        def:normal}
    Let $A$ and $B$ have the same physical alphabet and positive target bond
    dimension $D_A>0$.  Suppose that
    $\ket{V^{(N)}(B)}=\ket{V^{(N)}(A)}$ for every $N>0$.
    If $A$ is injective, then there exist
    $V\in\C^{D_A\times D_B}$ and $W\in\C^{D_B\times D_A}$ forming a
    reduction from $B$ to $A$.  The same conclusion holds when $A$ is normal.

    This is the $\lambda=1$ equality case of
    \cite[Proposition~20, lines~3815--3938]{Molnar2018SemiInjective}.  The unrestricted
    proportional premise does not imply the unscaled conclusion
    \cite{gap:mgsc18_reduction_proportionality_scalar}.
\end{theorem}

\begin{proof}
    \leanok
    \uses{lem:one_block_injective_of_injective,
        lem:isNBlkInjective_iff_blockTensor_isInjective,
        thm:same_mpv2_pos_block_tensor,
        thm:mps_reduction_cross_trace_open,
        thm:mpu_semisimple_scaled_trace_rank_one,
        thm:rank_one_factorizations}
    First suppose that $A$ is injective; equivalently, use blocking length one.
    More generally, choose any positive blocking length $p$ for which the
    blocked target $\widetilde A=A^{[p]}$ is injective, and put
    $\widetilde B=B^{[p]}$.  Positive blocking preserves equality of the MPV
    families.  Let $C$ be the coefficient dual of $\widetilde A$ and set
    \begin{align}
        K
         & =\sum_i \widetilde B^i\otimes(C^i)^{\mathsf T}.
        \notag
    \end{align}
    Regard $K$ as an endomorphism and take its Jordan--Chevalley decomposition
    $K=N+S$, where $N$ is nilpotent, $S$ is semisimple, and both belong to the
    algebra generated by $K$.  Their membership in this commutative generated
    algebra implies $NS=SN$.  The commuting nilpotent part does not change
    traces of powers.  Since $\tr K^n=D_A^n$ for every $n>0$, the semisimple
    endomorphism $D_A^{-1}S$ has trace one at every power above one.  The
    semisimple trace-power theorem gives $\rank S=1$; commutation and
    nilpotence then imply $NS=SN=0$ and
    \begin{align}
        K^J
         & =S^J
        \notag
    \end{align}
    for some $J>0$.

    Factor the matrix of $S$ with the mixed-bond orientation
    \begin{align}
        S_{(b,a),(b',a')}
         & =W_{ba}V_{a'b'}.
        \notag
    \end{align}
    Its first trace is $D_A$, so multiplication of outer products gives
    \begin{align}
        S^J
         & =D_A^{J-1}S.
        \notag
    \end{align}
    Apply the padded open-boundary identity with an arbitrary original-scale
    word $w$.  Rewriting its cross sum with the preceding factorization yields
    \begin{align}
        D_A^{J-1}(VB^wW)_{a'a}
         & =D_A^{J-1}(A^w)_{a'a}.
        \notag
    \end{align}
    Since $D_A>0$, cancel the nonzero scalar and conclude
    $VB^wW=A^w$ for every word $w$.  The empty word gives $VW=1_{D_A}$.
    For normal $A$, choose the positive injective blocking supplied by
    normality and repeat the same padded argument; the lower tail remains the
    unblocked word $w$, so the conclusion is already at the original physical
    scale.
\end{proof}

\begin{lemma}[Scalar covariance of rectangular reductions]
    \label{lem:mps_reduction_scalar_covariance}
    \lean{MPSTensor.IsReduction.smul,
        MPSTensor.IsReduction.evalWord_smul_target}
    \leanok
    \uses{def:mps_rectangular_reduction}
    If $V,W$ form a reduction from $B$ to $A$, then they also form a reduction
    from $cB$ to $cA$ for every $c\in\C$.  If instead they form a reduction
    from $B$ to $cA$, then
    \begin{align}
        VW & =1_{D_A},
           & VB^wW     & =c^{|w|}A^w
        \qquad\text{for every word }w.
        \label{eq:mps_reduction_scalar_covariance}
    \end{align}
\end{lemma}

\begin{proof}
    \leanok
    \uses{lem:eval_word_smul}
    Scaling both tensors multiplies each length-$|w|$ word by the same factor
    $c^{|w|}$, so the reduction identities are preserved.  For a reduction to
    $cA$, expand $(cA)^w=c^{|w|}A^w$ in the defining word identity.  The
    retraction $VW=1_{D_A}$ is unchanged.
\end{proof}

\begin{theorem}[Nonzero proportional existence of rectangular reductions]
    \label{thm:mps_proportional_reduction_exists}
    \lean{MPSTensor.exists_isReduction_smul_of_isInjective_of_mpv_eq_pow_mul,
        MPSTensor.exists_isReduction_smul_of_isNormal_of_mpv_eq_pow_mul}
    \leanok
    \uses{def:mpv, def:injective, def:normal,
        def:mps_rectangular_reduction}
    Let $A$ and $B$ have the same physical alphabet and positive target bond
    dimension $D_A>0$.  Let $c\in\C$ be nonzero, and suppose that
    \begin{align}
        V^{(N)}(B)_\sigma
         & =c^N V^{(N)}(A)_\sigma
        \qquad\text{for every }N>0\text{ and every configuration }\sigma.
        \label{eq:mps_proportional_reduction_hypothesis}
    \end{align}
    If $A$ is injective, then there exist
    $V\in\C^{D_A\times D_B}$ and $W\in\C^{D_B\times D_A}$ such that
    \begin{align}
        VW & =1_{D_A},
           & VB^wW     & =c^{|w|}A^w
        \qquad\text{for every word }w.
        \label{eq:mps_proportional_reduction_word_relation}
    \end{align}
    The same conclusion holds when $A$ is normal.

    This is the corrected nonzero-scalar form of
    \cite[Proposition~20]{Molnar2018SemiInjective}.  The unscaled conclusion
    printed there is false for unrestricted $c\ne1$, and the case $c=0$ does
    not admit this rescaling argument
    \cite{gap:mgsc18_reduction_proportionality_scalar}.
\end{theorem}

\begin{proof}
    \leanok
    \uses{lem:eval_word_smul, thm:mps_equality_reduction_exists,
        lem:mps_reduction_scalar_covariance}
    Set $\widetilde B=c^{-1}B$.  Expand its word evaluations with
    Lemma~\ref{lem:eval_word_smul}, take the trace, and use
    (\ref{eq:mps_proportional_reduction_hypothesis}).  The positive-length MPV
    families of $\widetilde B$ and $A$ are equal.  Apply
    Theorem~\ref{thm:mps_equality_reduction_exists} to obtain a reduction from
    $\widetilde B$ to $A$.  Scale both tensors by $c$ to obtain a reduction
    from $B$ to $cA$, and then apply
    Lemma~\ref{lem:mps_reduction_scalar_covariance} to get
    (\ref{eq:mps_proportional_reduction_word_relation}).
\end{proof}

\begin{lemma}[Nonempty products in a generated subsemigroup]
    \label{lem:subsemigroup_closure_nonempty_product}
    \lean{Subsemigroup.exists_nonempty_list_prod_of_mem_closure}
    \leanok
    Let $M$ be a monoid and let $S\subseteq M$.  Every element
    $x\in\langle S\rangle_{\mathrm{sg}}$ of the subsemigroup generated by
    $S$ admits a factorization
    \begin{align}
        x & =a_1\cdots a_n,
          & n               & \geq 1,
          & a_j             & \in S\quad(1\leq j\leq n).
        \notag
    \end{align}
\end{lemma}

\begin{proof}
    \leanok
    A generator has the required one-factor expression.  The product of two
    such expressions is represented by their concatenation, which is again
    nonempty and has every factor in $S$.
\end{proof}


\begin{definition}[Residual algebra of a selected reduction]
    \label{def:mps_reduction_residual_algebra}
    \lean{MPSTensor.reductionResidual,
        MPSTensor.reductionResidualGeneratorSet,
        MPSTensor.reductionResidualAlgebra}
    \leanok
    \uses{def:mps_rectangular_reduction}
    Let $(V,W)$ be a reduction from $B$ to $A$.  For each physical letter
    $i$, define
    \begin{align}
        N^i & = B^i-WA^iV.
        \label{eq:mps_reduction_residual_letter}
    \end{align}
    The \emph{residual algebra} $\mathcal R(B,A;V,W)$ is the nonunital
    subalgebra of $\MN{D_B}$ generated by the matrices $N^i$.
    This is the residual algebra of Proposition~21 in
    \cite{Molnar2018SemiInjective}.
\end{definition}

\begin{theorem}[Residual letters belong to the residual algebra]
    \label{thm:mps_reduction_residual_mem_algebra}
    \lean{MPSTensor.reductionResidual_mem_reductionResidualAlgebra}
    \leanok
    \uses{def:mps_reduction_residual_algebra}
    For every physical letter $i$, one has
    $N^i\in\mathcal R(B,A;V,W)$.
\end{theorem}

\begin{proof}
    \leanok
    Each residual letter is one of the generators of the residual algebra.
\end{proof}

\begin{definition}[Residual nilpotency length]
    \label{def:mps_reduction_residual_nilpotency_length}
    \lean{MPSTensor.IsReductionResidualNilpotencyBound,
        MPSTensor.reductionResidualNilpotencyLength}
    \leanok
    \uses{def:mps_reduction_residual_algebra}
    A natural number $L$ is a \emph{residual-word nilpotency bound} if
    every word $N^{i_1}\cdots N^{i_L}$ of exactly $L$ residual letters is
    zero.  The residual nilpotency length $\ell(B,A;V,W)$ is the infimum
    of the set of exact-length bounds.  Definition~8 of
    \cite{Molnar2018SemiInjective} instead asks that every residual word of
    every length $n\geq L$ vanish; the passage from an exact-length bound to
    this consequence is stated separately in
    Theorem~\ref{thm:mps_reduction_residual_bound_consequences}.
\end{definition}

\begin{theorem}[Exterior-word strengthening of the residual sandwich]
    \label{thm:mps_reduction_residual_exterior_sandwich}
    \lean{MPSTensor.IsReduction.evalWord_mul_reductionResidual_mul_evalWord_eq_zero}
    \leanok
    \uses{def:mps_rectangular_reduction,
        def:mps_reduction_residual_algebra}
    Let $(V,W)$ be a reduction from $B$ to $A$.  For arbitrary exterior words
    $\mathbf p,\mathbf q$ and every nonempty word $\mathbf u$,
    \begin{align}
        VB^{\mathbf p}N^{\mathbf u}B^{\mathbf q}W & =0.
        \label{eq:mps_reduction_residual_exterior_sandwich}
    \end{align}
    This strengthens the sandwiched residual identity below by allowing
    arbitrary exterior $B$-words.
\end{theorem}

\begin{proof}
    \leanok
    \uses{thm:mps_rectangular_reduction_identities}
    Apply $VB^{\mathbf w}W=A^{\mathbf w}$ to $\mathbf p$, $\mathbf q$, and
    $\mathbf p\,i\,\mathbf q$.  The two expressions for
    $A^{\mathbf p}A^iA^{\mathbf q}$ give
    \begin{align}
        VB^{\mathbf p}N^iB^{\mathbf q}W & =0.
        \notag
    \end{align}
    Induction on the nonempty residual word $\mathbf u$ proves
    (\ref{eq:mps_reduction_residual_exterior_sandwich}).
\end{proof}

\begin{theorem}[Positive residual words vanish under the reduction]
    \label{thm:mps_reduction_residual_sandwich}
    \lean{MPSTensor.IsReduction.evalWord_reductionResidual_sandwich_eq_zero}
    \leanok
    \uses{def:mps_rectangular_reduction,
        def:mps_reduction_residual_algebra}
    Let $(V,W)$ be a reduction from $B$ to $A$.  For every nonempty word
    $\mathbf u$,
    \begin{align}
        VN^{\mathbf u}W & =0
        \qquad (|\mathbf u|>0).
        \label{eq:mps_reduction_residual_sandwich}
    \end{align}
    No equality of closed chains is required.
    This is the sandwiched residual identity of
    \cite{Molnar2018SemiInjective}.
\end{theorem}

\begin{proof}
    \leanok
    \uses{thm:mps_reduction_residual_exterior_sandwich}
    Take both exterior words empty in
    (\ref{eq:mps_reduction_residual_exterior_sandwich}).
\end{proof}

\begin{definition}[Strict residual-window sums]
    \label{def:mps_reduction_strict_window_sums}
    \lean{MPSTensor.reductionInitialWindowSum,
        MPSTensor.reductionStrictWindowSum,
        MPSTensor.reductionInitialWindowRangeSum,
        MPSTensor.reductionStrictWindowRangeSum}
    \leanok
    \uses{def:mps_rectangular_reduction,
        def:mps_reduction_residual_algebra, def:eval_word}
    Let $(V,W)$ be a reduction from $B$ to $A$.  For a word $\mathbf w$, set
    $\mathcal I(\varnothing)=\mathcal W(\varnothing)=0$ and define
    recursively
    \begin{align}
        \mathcal I(i\mathbf w)
         & =WA^iV\bigl(N^{\mathbf w}+\mathcal I(\mathbf w)\bigr),
        \label{eq:mps_reduction_initial_window_recursion}         \\
        \mathcal W(i\mathbf w)
         & =\mathcal I(i\mathbf w)+N^i\mathcal W(\mathbf w).
        \label{eq:mps_reduction_strict_window_recursion}
    \end{align}
    For $\mathbf i=(i_1,\ldots,i_n)$, also set
    \begin{align}
        \widehat{\mathcal I}(\mathbf i)
         & =\sum_{m=1}^{n}
        W A^{i_1}\cdots A^{i_m}V
        N^{i_{m+1}}\cdots N^{i_n},           \\
        \widehat{\mathcal W}(\mathbf i)
         & =\sum_{r=0}^{n-1}\sum_{m=1}^{n-r}
        N^{i_1}\cdots N^{i_r}
        W A^{i_{r+1}}\cdots A^{i_{r+m}}V
        N^{i_{r+m+1}}\cdots N^{i_n}.
        \label{eq:mps_reduction_strict_window_sum}
    \end{align}
    The substitution $s=r+m$ identifies the second indexing set with
    $0\leq r<s\leq n$.  Thus the central $A$-word is nonempty, and
    $\widehat{\mathcal I}$ is the part with $r=0$.
\end{definition}

\begin{theorem}[Recursive residual expansion]
    \label{thm:mps_reduction_recursive_residual_expansion}
    \lean{MPSTensor.IsReduction.evalWord_eq_reductionResidual_add_strictWindowSum}
    \leanok
    \uses{def:mps_rectangular_reduction,
        def:mps_reduction_strict_window_sums}
    Under the reduction hypothesis, every word satisfies
    \begin{align}
        B^{\mathbf i}=N^{\mathbf i}+\mathcal W(\mathbf i).
        \label{eq:mps_reduction_recursive_residual_expansion}
    \end{align}
    Combining this recursion with the explicit window sum below gives the
    corrected residual expansion.
\end{theorem}

\begin{proof}
    \leanok
    \uses{thm:mps_reduction_residual_sandwich,
        thm:mps_rectangular_reduction_identities}
    Induct on the word and split the first letter as $B^i=N^i+WA^iV$.
    The induction hypothesis supplies the recursive windows in the remaining
    word.  By (\ref{eq:mps_reduction_residual_sandwich}),
    $VN^{\mathbf u}W=0$ for every nonempty $\mathbf u$, so prefixing by $V$
    collapses $\mathcal W(\mathbf w)$ to $\mathcal I(\mathbf w)$; adjacent
    selected blocks combine by $VW=1$.
\end{proof}

\begin{theorem}[All-cut contracted residual identities]
    \label{thm:mps_reduction_all_cut_contracted}
    \lean{MPSTensor.reductionLeftContractedAllCutSum,
        MPSTensor.reductionRightContractedAllCutSum,
        MPSTensor.reductionLeftContractedSum,
        MPSTensor.reductionRightContractedSum,
        MPSTensor.reductionLeftContractedSum_eq_allCutSum,
        MPSTensor.reductionRightContractedSum_eq_allCutSum,
        MPSTensor.IsReduction.mul_evalWord_eq_reductionLeftContractedSum,
        MPSTensor.IsReduction.evalWord_mul_eq_reductionRightContractedSum}
    \leanok
    \uses{def:mps_rectangular_reduction,
        def:mps_reduction_residual_algebra}
    Let $(V,W)$ be a reduction from $B$ to $A$, and write
    $\mathbf i=(i_1,\ldots,i_n)$.  Then
    \begin{align}
        VB^{i_1}\cdots B^{i_n}
         & =\sum_{s=0}^{n}
        A^{i_1}\cdots A^{i_s}V
        N^{i_{s+1}}\cdots N^{i_n},
        \label{eq:mps_reduction_left_all_cut} \\
        B^{i_1}\cdots B^{i_n}W
         & =\sum_{r=0}^{n}
        N^{i_1}\cdots N^{i_r}W
        A^{i_{r+1}}\cdots A^{i_n}.
        \label{eq:mps_reduction_right_all_cut}
    \end{align}
    These identities require only the reduction.  They are precursors to the
    bound-truncated formulas in
    Theorem~\ref{thm:mps_reduction_contracted_expansions}.
\end{theorem}

\begin{proof}
    \leanok
    \uses{thm:mps_rectangular_reduction_identities,
        thm:mps_reduction_residual_sandwich,
        thm:mps_reduction_recursive_residual_expansion}
    Prove both identities by induction on the word.  For the left contraction,
    decompose the first letter and apply
    (\ref{eq:mps_reduction_recursive_residual_expansion}) to the tail; the
    sandwiched terms $VN^{\mathbf u}W$ vanish by
    (\ref{eq:mps_reduction_residual_sandwich}), and $VW=1$ merges adjacent
    selected blocks.  For the right contraction, decompose the first letter
    and apply the reduction identity to the tail.
\end{proof}

\begin{theorem}[Recursive and explicit strict-window sums]
    \label{thm:mps_reduction_strict_window_range}
    \lean{MPSTensor.IsReduction.reductionStrictWindowSum_eq_reductionStrictWindowRangeSum}
    \leanok
    \uses{def:mps_rectangular_reduction,
        def:mps_reduction_strict_window_sums}
    Under the reduction hypothesis, for every word $\mathbf i$,
    \begin{align}
        \mathcal W(\mathbf i)=\widehat{\mathcal W}(\mathbf i).
    \end{align}
\end{theorem}

\begin{proof}
    \leanok
    \uses{thm:mps_reduction_all_cut_contracted,
        thm:mps_reduction_recursive_residual_expansion,
        thm:mps_reduction_residual_sandwich}
    Induct on the length of the word.  The summands with $r=0$ constitute
    $\widehat{\mathcal I}$.  By the left all-cut identity,
    \begin{align}
        \widehat{\mathcal I}(i_1,\ldots,i_n)
        =WA^{i_1}VB^{i_2}\cdots B^{i_n}
        =\mathcal I(i_1,\ldots,i_n),
    \end{align}
    where the last equality is the initial-window case of the recursive
    expansion.  Deleting the first letter identifies the summands with $r>0$
    with $N^{i_1}\widehat{\mathcal W}(i_2,\ldots,i_n)$.
\end{proof}

\begin{theorem}[Corrected strict-window residual expansion]
    \label{thm:mps_reduction_residual_expansion}
    \lean{MPSTensor.IsReduction.evalWord_eq_reductionResidual_add_strictWindowRangeSum}
    \leanok
    \uses{def:mps_rectangular_reduction,
        def:mps_reduction_strict_window_sums}
    Let $(V,W)$ be a reduction from $B$ to $A$.  For every word
    $\mathbf i=(i_1,\ldots,i_n)$,
    \begin{align}
        B^{i_1}\cdots B^{i_n}
         & =N^{i_1}\cdots N^{i_n}
        +\sum_{0\leq r<s\leq n}
        N^{i_1}\cdots N^{i_r}
        W A^{i_{r+1}}\cdots A^{i_s}V
        N^{i_{s+1}}\cdots N^{i_n}.
        \label{eq:mps_reduction_residual_expansion}
    \end{align}
    This statement requires only the reduction. The range $0\leq r<s\leq n$
    ensures nonempty central words; see
    \cite{gap:mgsc18_residual_expansion_empty_word_error}.
\end{theorem}

\begin{proof}
    \leanok
    \uses{thm:mps_reduction_recursive_residual_expansion,
        thm:mps_reduction_strict_window_range}
    Substitute
    $\mathcal W(\mathbf i)=\widehat{\mathcal W}(\mathbf i)$ into
    (\ref{eq:mps_reduction_recursive_residual_expansion}), and then reindex
    the positive central length by $s=r+m$.
\end{proof}

\begin{theorem}[Residual trace vanishing and elementwise nilpotence]
    \label{thm:mps_reduction_residual_algebra_nilpotent}
    \lean{MPSTensor.IsReduction.trace_evalWord_eq_trace_evalWord_add_trace_reductionResidual,
        MPSTensor.IsReduction.trace_evalWord_reductionResidual_eq_zero,
        MPSTensor.IsReduction.trace_eq_zero_of_mem_reductionResidualAlgebra,
        MPSTensor.IsReduction.isNilpotent_of_mem_reductionResidualAlgebra}
    \leanok
    \uses{def:same_mpv2_pos, def:mps_rectangular_reduction,
        def:mps_reduction_residual_algebra}
    Let $(V,W)$ be a reduction from $B$ to $A$, and assume
    $\mc{V}^+(B)=\mc{V}^+(A)$.  Then every nonempty residual word has trace
    zero, every element of $\mathcal R(B,A;V,W)$ has trace zero, and every
    element of this algebra is nilpotent.

    This establishes the nil-algebra step of
    \cite[Proposition~21]{Molnar2018SemiInjective} in the equality case
    $\lambda=1$; see \cite{gap:mgsc18_reduction_proportionality_scalar}.
\end{theorem}

\begin{proof}
    \leanok
    \uses{thm:mps_reduction_residual_expansion,
        thm:mps_reduction_residual_exterior_sandwich,
        thm:mps_reduction_residual_sandwich,
        thm:mps_reduction_strict_window_range,
        thm:mps_rectangular_reduction_identities,
        thm:mps_reduction_cross_trace_open,
        lem:subsemigroup_closure_nonempty_product,
        thm:mpu_shifted_zero_trace_power_nilpotent}
    Take the trace of the corrected residual expansion.  For a mixed term in
    the explicit strict-window sum, cyclicity gives
    \begin{align}
        \tr\!\left(
        N^{\mathbf p}WA^{\mathbf a}VN^{\mathbf q}\right)
         & =\tr\!\left(
        \left(VN^{\mathbf q}N^{\mathbf p}W\right)A^{\mathbf a}\right)
        =0,
        \notag
    \end{align}
    by (\ref{eq:mps_reduction_residual_sandwich}).  In the recursive sum, the
    corresponding step is
    \begin{align}
        \tr\!\left(N^{\mathbf u}WA^iVB^{\mathbf w}\right)
         & =\tr\!\left(
        \left(VB^{\mathbf w}N^{\mathbf u}W\right)A^i\right)
        =0,
        \notag
    \end{align}
    by (\ref{eq:mps_reduction_residual_exterior_sandwich}), for every
    nonempty residual word $\mathbf u$.  The unique term with no residual
    letters is $WA^{\mathbf i}V$, whose trace is $\tr(A^{\mathbf i})$ because
    $VW=1$.  Thus
    \begin{align}
        \tr(B^{\mathbf i})
         & =\tr(A^{\mathbf i})+\tr(N^{\mathbf i}).
        \label{eq:mps_reduction_residual_trace_decomposition}
    \end{align}
    Equality of the two closed chains now gives
    $\tr(N^{\mathbf i})=0$.  Linearity extends this identity to the generated
    nonunital algebra.  Vanishing traces of positive powers make each algebra
    element nilpotent.
\end{proof}

\begin{theorem}[Dimension bound for the residual nilpotency length]
    \label{thm:mps_reduction_residual_dimension_bound}
    \lean{MPSTensor.IsReduction.reductionResidualAlgebra_list_prod_eq_zero,
        MPSTensor.IsReduction.bondDim_isReductionResidualNilpotencyBound,
        MPSTensor.IsReduction.reductionResidualNilpotencyLength_isBound,
        MPSTensor.IsReduction.reductionResidualNilpotencyLength_le_bondDim,
        MPSTensor.IsReduction.evalWord_reductionResidual_eq_zero_of_bondDim_le_length}
    \leanok
    \uses{def:same_mpv2_pos, def:mps_rectangular_reduction,
        def:mps_reduction_residual_algebra,
        def:mps_reduction_residual_nilpotency_length}
    Let $(V,W)$ be a reduction from $B$ to $A$, and assume
    $\mc{V}^+(B)=\mc{V}^+(A)$.  Then every mixed product of $D_B$ elements of
    $\mathcal R(B,A;V,W)$ vanishes.  Consequently $D_B$ is a residual-word
    nilpotency bound.  The least bound exists, is itself a bound, and satisfies
    \begin{align}
        \ell(B,A;V,W) & \leq D_B.
        \label{eq:mps_reduction_residual_nilpotency_length_bound}
    \end{align}
    In particular, $N^{\mathbf i}=0$ whenever $|\mathbf i|\geq D_B$.
    The resulting vanishing of all sufficiently long mixed products is the
    nilpotency conclusion of Proposition~21 in
    \cite{Molnar2018SemiInjective}.  The specific numerical estimate by $D_B$
    follows from the finite-dimensional nil-matrix theorem.
\end{theorem}

\begin{proof}
    \leanok
    \uses{thm:mps_reduction_residual_algebra_nilpotent,
        thm:mps_reduction_residual_mem_algebra,
        thm:mps_reduction_residual_bound_consequences,
        thm:mpu_nil_matrices}
    Apply the finite-dimensional nil-matrix theorem to the residual algebra.
    Its dimension-sharp matrix form makes every product of $D_B$ algebra
    elements zero.  Residual letters lie in this algebra, so $D_B$ is a bound.
    Well-ordering gives the least bound, and prefix factorization gives
    vanishing for all longer residual words.
\end{proof}

\begin{theorem}[Consequences of an arbitrary residual bound]
    \label{thm:mps_reduction_residual_bound_consequences}
    \lean{MPSTensor.IsReduction.evalWord_reductionResidual_eq_zero_of_bound_le_length}
    \leanok
    \uses{def:mps_reduction_residual_nilpotency_length}
    If $L$ is a residual-word nilpotency bound and
    $|\mathbf i|\geq L$, then
    \begin{align}
        N^{\mathbf i} & =0.
        \label{eq:mps_reduction_residual_bound_word}
    \end{align}
    This conclusion uses only the assumed bound; it does not require equality
    of closed chains.
\end{theorem}

\begin{proof}
    \leanok
    Split the residual word after its first $L$ letters.  The prefix is zero
    by the bound, so the full word is zero.
\end{proof}

\begin{theorem}[Bound-truncated contracted residual expansions]
    \label{thm:mps_reduction_contracted_expansions}
    \lean{MPSTensor.reductionLeftContractedRangeSum,
        MPSTensor.reductionRightContractedRangeSum,
        MPSTensor.IsReduction.mul_evalWord_eq_reductionLeftContractedRangeSum,
        MPSTensor.IsReduction.evalWord_mul_eq_reductionRightContractedRangeSum}
    \leanok
    \uses{def:mps_rectangular_reduction,
        def:mps_reduction_residual_nilpotency_length}
    Let $(V,W)$ be a reduction from $B$ to $A$, let $L$ be any residual-word
    nilpotency bound, and write
    $\mathbf i=(i_1,\ldots,i_n)$.  Then
    \begin{align}
        VB^{i_1}\cdots B^{i_n}
         & =\sum_{n-L\leq s\leq n}
        A^{i_1}\cdots A^{i_s}V
        N^{i_{s+1}}\cdots N^{i_n},
        \label{eq:mps_reduction_left_contracted_expansion} \\
        B^{i_1}\cdots B^{i_n}W
         & =\sum_{0\leq r\leq\min(L,n)}
        N^{i_1}\cdots N^{i_r}W
        A^{i_{r+1}}\cdots A^{i_n}.
        \label{eq:mps_reduction_right_contracted_expansion}
    \end{align}
    Here $n-L$ denotes truncated subtraction, equivalently
    $\max(0,n-L)$ in integer notation.  The endpoints with a residual word of
    length exactly $L$ are retained as zero terms.  These are the two
    contracted formulas of \cite{Molnar2018SemiInjective}.
\end{theorem}

\begin{proof}
    \leanok
    \uses{thm:mps_reduction_all_cut_contracted,
        thm:mps_reduction_residual_bound_consequences}
    Start from (\ref{eq:mps_reduction_left_all_cut}) and
    (\ref{eq:mps_reduction_right_all_cut}).  In the left sum, every omitted
    term has residual suffix length at least $L$.  In the right sum, every
    omitted term has residual prefix length at least $L$, so
    (\ref{eq:mps_reduction_residual_bound_word}) makes all these terms zero.
\end{proof}

\begin{theorem}[Exterior-buffer identity]
    \label{thm:mps_reduction_exterior_buffer}
    \lean{MPSTensor.IsReduction.evalWord_mul_reduced_exterior_eq_evalWord_append}
    \leanok
    \uses{def:mps_rectangular_reduction,
        def:mps_reduction_residual_nilpotency_length}
    Let $(V,W)$ be a reduction from $B$ to $A$, and let $L$ be any
    residual-word nilpotency bound.  If the central word
    $\mathbf c$ is nonempty and the exterior words $\mathbf p,\mathbf q$
    satisfy
    \begin{align}
        L & \leq |\mathbf p|+1,
          & L                   & \leq |\mathbf q|+1,
        \label{eq:mps_reduction_exterior_buffer_bounds}
    \end{align}
    then
    \begin{align}
        B^{\mathbf p}WA^{\mathbf c}VB^{\mathbf q}
         & =B^{\mathbf p\mathbf c\mathbf q}.
        \label{eq:mps_reduction_exterior_buffer}
    \end{align}
    No closed-chain equality is assumed: only the reduction and the supplied
    residual bound enter this identity.  It follows from the corrected
    residual expansion by eliminating every term with a sufficiently long
    residual buffer.
\end{theorem}

\begin{proof}
    \leanok
    \uses{thm:mps_rectangular_reduction_identities,
        thm:mps_reduction_residual_exterior_sandwich,
        thm:mps_reduction_recursive_residual_expansion,
        thm:mps_reduction_residual_bound_consequences}
    First use induction along the right exterior word to show that a positive
    residual block preceded by $V$ vanishes once its length together with the
    right buffer reaches $L$.  Expand only the left exterior word
    $B^{\mathbf p}$ by
    Theorem~\ref{thm:mps_reduction_recursive_residual_expansion}.  The left
    buffer bound removes the pure residual term, while the preceding induction
    and (\ref{eq:mps_reduction_residual_exterior_sandwich}) remove the
    strict-window terms.
    Hence
    $B^{\mathbf p}N^iB^{\mathbf q}=0$ under the two buffer bounds.

    Separately, induction on a word $\mathbf a$ gives
    $VB^{\mathbf a\mathbf q}=A^{\mathbf a}VB^{\mathbf q}$ from the right
    buffer bound.  Write the nonempty central word as
    $\mathbf c=i\mathbf a$ and decompose the corresponding letter as
    $B^i=N^i+WA^iV$.  The residual term is zero by the first step, and the
    second induction identifies the remaining term with
    (\ref{eq:mps_reduction_exterior_buffer}).
\end{proof}

\begin{definition}[Exterior buffer length of a reduction]
    \label{def:mps_reduction_exterior_buffer_length}
    \lean{MPSTensor.IsReductionExteriorBufferLength,
        MPSTensor.IsReductionExteriorBufferLength.mono}
    \leanok
    \uses{def:mps_rectangular_reduction}
    Let $(V,W)$ be a reduction from $B$ to $A$.  A natural number $m$ is an
    \emph{exterior buffer length} for $(V,W)$ if
    \begin{align}
        B^{\mathbf p}WA^{\mathbf c}VB^{\mathbf q}
         & =B^{\mathbf p\mathbf c\mathbf q}
        \label{eq:mps_reduction_exterior_buffer_length}
    \end{align}
    for every nonempty word $\mathbf c$ and all words $\mathbf p,\mathbf q$
    with $|\mathbf p|\geq m$ and $|\mathbf q|\geq m$.  Every $m'\geq m$ is
    then also an exterior buffer length.  This is the exterior fusion
    equation of \cite[lines~1409--1498]{FrancoRubio2025MPUGauging}, with
    exterior length $m\geq\ell$ and central length $|\mathbf c|=n+1$, $n\geq0$.
\end{definition}

\begin{theorem}[Reductions and exterior buffers under blocking]
    \label{thm:mps_reduction_blocking}
    \lean{MPSTensor.IsReduction.isReductionExteriorBufferLength_of_bound,
        MPSTensor.IsReduction.isReductionExteriorBufferLength_nilpotencyLength,
        MPSTensor.IsReduction.blockTensor,
        MPSTensor.IsReduction.reindexPhysical,
        MPSTensor.IsReductionExteriorBufferLength.blockTensor,
        MPSTensor.IsReductionExteriorBufferLength.reindexPhysical}
    \leanok
    \uses{def:mps_rectangular_reduction,
        def:mps_reduction_exterior_buffer_length,
        def:mps_reduction_residual_nilpotency_length,
        def:same_mpv2_pos, def:blocked_tensor}
    Let $(V,W)$ be a reduction from $B$ to $A$.  If $N$ is a residual-word
    nilpotency bound, then $N-1$ is an exterior buffer length; in particular,
    if $\mc{V}^+(B)=\mc{V}^+(A)$, then $\ell(B,A;V,W)-1$ is an exterior
    buffer length.

    For every $L$, the same pair $(V,W)$ is a reduction from $B^{[L]}$ to
    $A^{[L]}$, and for every map $f$ between physical alphabets it is a
    reduction from $B\circ f$ to $A\circ f$.  If $m$ is an exterior buffer
    length for $(V,W)$, $L>0$, and $m\leq kL$, then $k$ is an exterior buffer
    length for $(V,W)$ as a reduction from $B^{[L]}$ to $A^{[L]}$; reindexing
    both tensors by the same map $f$ also preserves exterior buffer lengths.
\end{theorem}

\begin{proof}
    \leanok
    \uses{thm:mps_reduction_exterior_buffer,
        thm:mps_reduction_residual_dimension_bound,
        lem:eval_word_block}
    The bounds $N\leq|\mathbf p|+1$ and $N\leq|\mathbf q|+1$ in
    Theorem~\ref{thm:mps_reduction_exterior_buffer} read
    $N-1\leq|\mathbf p|$ and $N-1\leq|\mathbf q|$, and
    Theorem~\ref{thm:mps_reduction_residual_dimension_bound} makes the least
    length a bound under equality of the positive-length closed chains.
    Word evaluation of $B^{[L]}$ on a blocked word is evaluation of $B$ on
    the flattened word, and evaluation of $B\circ f$ on a word is evaluation
    of $B$ on the mapped word, so the reduction identities transfer.  For the
    exterior identity, the flattened exterior words have length at least
    $kL\geq m$ and the flattened central word has positive length $|\mathbf c|L$.
\end{proof}

\begin{definition}[Common buffer length of a finite family]
    \label{def:common_buffer_length}
    \lean{commonBufferLength, commonBufferLength_pos,
        le_commonBufferLength_add_one, commonBufferLength_le}
    \leanok
    For a finite family of natural numbers $(N_i)_{i\in I}$, the
    \emph{common buffer length} is
    \begin{align}
        L & =\max\Bigl(1,\max_{i\in I}(N_i-1)\Bigr),
        \label{eq:common_buffer_length}
    \end{align}
    with $N_i-1$ read as $0$ when $N_i=0$.  It is positive, satisfies
    $N_i\leq L+1$ for every $i\in I$, and is the least positive integer with
    this property.
\end{definition}

\begin{theorem}[Common blocking of a finite family of reductions]
    \label{thm:mps_reduction_common_blocking}
    \lean{MPSTensor.IsReduction.blockTensor_commonBufferLength_one}
    \leanok
    \uses{def:mps_rectangular_reduction,
        def:mps_reduction_exterior_buffer_length,
        def:mps_reduction_residual_nilpotency_length,
        def:common_buffer_length, def:blocked_tensor}
    Let $(N_i)_{i\in I}$ be a finite family of natural numbers with common
    buffer length $L$, and let $(V,W)$ be a reduction from $B$ to $A$ for
    which $N_i$ is a residual-word nilpotency bound for some $i\in I$.  Then
    $1$ is an exterior buffer length for $(V,W)$ as a reduction from
    $B^{[L]}$ to $A^{[L]}$: (\ref{eq:mps_reduction_exterior_buffer_length})
    holds for the blocked tensors whenever the blocked words
    $\mathbf p,\mathbf c,\mathbf q$ are nonempty.

    This formalizes the reduction of the buffer length to $1$ by blocking in
    \cite{FrancoRubio2025MPUGauging} and
    \cite[Theorem~1 footnote]{GarreRubioSchuch2024DomainWall}; see
    \cite{gap:mgsc18_nilpotency_length_one_terminology}.
\end{theorem}

\begin{proof}
    \leanok
    \uses{thm:mps_reduction_blocking}
    By Theorem~\ref{thm:mps_reduction_blocking}, $N_i-1$ is an exterior
    buffer length, and $N_i-1\leq L=1\cdot L$ with $L>0$ by
    (\ref{eq:common_buffer_length}).  The blocking statement of the same
    theorem with $k=1$ gives the conclusion.
\end{proof}

\begin{theorem}[Boundary-dressed uniqueness of rectangular reductions]
    \label{thm:mps_reduction_boundary_dressed_uniqueness}
    \lean{MPSTensor.IsReduction.exists_boundary_dressed_proportional,
        MPSTensor.IsReduction.exists_boundary_dressed_proportional_of_nilpotencyLength_le}
    \leanok
    \discussion{7395}
    \uses{def:same_mpv2_pos, def:mps_rectangular_reduction,
        def:mps_reduction_residual_nilpotency_length, def:normal}
    Let $A$ be normal, let $B$ have the same physical alphabet, and let
    $(V,W)$ and $(\widetilde V,\widetilde W)$ be two reductions from $B$ to
    $A$.  Suppose that $N_0\in\mathbb N$ is a residual-word nilpotency bound
    for both reductions: every word of length $N_0$ in either residual tensor
    vanishes.  Then there exists a single scalar $\lambda\in\C$ with
    $\lambda\ne0$ such
    that, for every word $\mathbf a$ satisfying
    \begin{align}
        |\mathbf a| & > 2N_0,
        \label{eq:mps_reduction_boundary_dressed_length}
    \end{align}
    one has
    \begin{align}
        VB^{\mathbf a}
         & =\lambda\,\widetilde V B^{\mathbf a},
        \label{eq:mps_reduction_boundary_dressed_left} \\
        B^{\mathbf a}W
         & =\lambda^{-1}B^{\mathbf a}\widetilde W.
        \label{eq:mps_reduction_boundary_dressed_right}
    \end{align}
    If $A$ and $B$ have the same positive-length closed-chain coefficients and
    the two least residual nilpotency lengths are at most $N_0$, then $N_0$ is
    such a common bound and the same conclusion follows.  The direct-bound
    statement itself requires no equality of closed-chain coefficients.

    Following \cite[Theorem~22]{Molnar2018SemiInjective}, these boundary-dressed
    identities hold for $|\mathbf a|>2N_0$.
\end{theorem}

\begin{proof}
    \leanok
    \uses{thm:mps_rectangular_reduction_identities,
        thm:mps_reduction_residual_dimension_bound,
        thm:mps_reduction_residual_bound_consequences,
        thm:mps_reduction_exterior_buffer, def:l_blk_injective,
        thm:exact_block_injectivity_positive_multiple}
    If the target bond dimension is zero, both dressed identities hold
    entrywise for $\lambda=1$ because the relevant target row or column index is
    empty.  Hence assume that the target bond dimension is positive.

    Choose $L>0$ such that $A$ is $L$-block injective.  For exterior words
    $\mathbf p,\mathbf q$ satisfying the $N_0$ buffer inequalities, compare the
    two exterior-buffer expansions of the mixed contraction
    $VB^{\mathbf p}\widetilde W$.  Inserting a nonempty central $A$-word gives
    \begin{align}
        (VB^{\mathbf p}\widetilde W)A^{\mathbf c}A^{\mathbf q}
        =A^{\mathbf p}A^{\mathbf c}(VB^{\mathbf q}\widetilde W).
        \notag
    \end{align}
    Since the length-$L$ words span the full target matrix algebra, linearity
    replaces $A^{\mathbf c}$ by an arbitrary matrix.  Insert matrix units and
    choose a nonzero target word at a positive multiple of $L$.  Entrywise
    cancellation then yields one scalar $\lambda$ with
    $VB^{\mathbf p}\widetilde W=\lambda A^{\mathbf p}$ for every sufficiently
    buffered $\mathbf p$.  Reversing the two reductions gives a scalar $\mu$
    for the opposite mixed boundary.

    Choose a longer nonzero target word, split it into two buffered exterior
    pieces and a nonempty injective central block, and apply the exterior-buffer
    identity once more.  The two mixed-boundary formulas give
    $(\lambda\mu)A^{\mathbf a}=A^{\mathbf a}$; nonvanishing of the chosen word
    implies $\lambda\mu=1$, so in particular $\lambda\ne0$.

    Finally, if $|\mathbf a|>2N_0$, split $\mathbf a$ into exterior words of
    length $N_0$ and a nonempty center.  Substitution of the two mixed-boundary
    formulas into the corresponding exterior-buffer identities gives
    (\ref{eq:mps_reduction_boundary_dressed_left}) and first gives
    $B^{\mathbf a}\widetilde W=\lambda B^{\mathbf a}W$ on the right.
    Multiplication by $\lambda^{-1}$ yields
    (\ref{eq:mps_reduction_boundary_dressed_right}).

    For the least-length corollary, positive-length closed-chain equality makes
    each least residual nilpotency length an actual bound.  Prefix
    factorization propagates vanishing from that least length to $N_0$, so the
    direct-bound theorem applies.
\end{proof}
%
